\documentclass[11pt]{article}
\usepackage[T1]{fontenc}
\usepackage[utf8]{inputenc}
\usepackage{lmodern}
\usepackage[margin=1in]{geometry}
\usepackage{amsmath,amssymb,amsthm,mathtools,bm}
\usepackage{booktabs,array,longtable}
\usepackage{enumitem}
\usepackage{hyperref}
\newcommand{\OPHPaperReleaseID}{r1339}
\newcommand{\OPHPaperReleaseDate}{March 14, 2026}
\newcommand{\OPHPaperReleaseBanner}{%
\begin{center}
\small\textbf{Paper release:} \texttt{\OPHPaperReleaseID}\quad\textbf{Released:} \OPHPaperReleaseDate
\end{center}
}


\hypersetup{colorlinks=true,linkcolor=blue,citecolor=blue,urlcolor=blue}
\setlength{\emergencystretch}{3em}
\setlength{\tabcolsep}{5pt}
\renewcommand{\arraystretch}{1.08}

\newtheorem{theorem}{Theorem}[section]
\newtheorem{corollary}[theorem]{Corollary}
\newtheorem{proposition}[theorem]{Proposition}
\newtheorem{lemma}[theorem]{Lemma}
\newtheorem{definition}[theorem]{Definition}
\newtheorem{axiom}[theorem]{Axiom}
\newtheorem{assumption}[theorem]{Assumption}
\newtheorem{remark}[theorem]{Remark}

\newcommand{\SU}{\mathrm{SU}}
\newcommand{\U}{\mathrm{U}}
\newcommand{\Conf}{\mathrm{Conf}}
\newcommand{\nf}{\operatorname{nf}}

\title{A Conditional Reconstruction Program for Relativity and Standard Model Structure from\\
Observer-Overlap Consistency\\[0.4em]
\large Observers Are All You Need (Compact)}
\author{B.\ M\"uller}
\date{March 13, 2026}

\begin{document}
\maketitle
\OPHPaperReleaseBanner

\begin{abstract}
We present a compact, falsifiable Phase-I statement of Observer-Patch Holography (OPH) as a conditional reconstruction program whose claimed recovered core is limited to relativity and Standard Model structure. From five axioms, the stated scaling-limit/categorical regularity package, and two external continuous inputs used only in the current quantitative implementation, the compact note derives a schedule-independent normal form for overlap consistency, a conditional Lorentz branch on the screen, a conditional Jacobson-type Einstein branch from fixed-cap generalized-entropy stationarity together with the null modular bridge and the scaling-limit null-stress premise, and conditional compact gauge reconstruction in the bosonic internal-gauge branch from a refinement-stable symmetric edge-sector colimit. After the extended Minimal Admissible Realization selector is imposed on the connected positive-dimensional Lie admissible class with one connected abelian charge factor, the realized low-energy gauge group on that realized branch is
\[
\frac{\SU(3)\times \SU(2)\times \U(1)}{\mathbb Z_6},
\qquad
N_c=3,
\qquad
N_g=3,
\]
together with the exact rational hypercharge lattice for the realized one-generation Standard Model chiral matter package with one Higgs doublet. The two external continuous inputs of the present implementation are the pixel area
\[
P \equiv a_{\mathrm{cell}}/\ell_P^2 = 1.63094
\]
and the total screen capacity
\[
N_{\mathrm{scr}} \equiv \log \dim \mathcal H_{\mathrm{tot}} \sim 10^{122}.
\]
In the current implementation, $P$ feeds the supplement-backed gauge/electroweak calibration sector and $N_{\mathrm{scr}}$ feeds the capacity branch for $\Lambda$; the D10--D11 numerical branches are reported only as implementation-level checks, not as the recovered core itself. Charged-lepton/Koide, texture, benchmark neutrino-model, dark-sector, baryogenesis, black-hole-spectroscopy, and string/worldsheet discussions are explicitly deferred continuations requiring additional ans\"atze beyond the recovered core and are not claimed as outputs of overlap consistency plus the stated scaling-limit and categorical premises. The purpose of this compact note is therefore to state a clean Phase-I claim set: relativity and Standard Model structure are the recovered core, while all downstream branches are either input-dependent or supplement-backed secondary sectors, or post-Phase-I phenomenology.
\end{abstract}

\section{Introduction}

The standard model of particle physics and general relativity are successful effective descriptions. They leave open a familiar list of questions: why the Lorentz group appears, why the low-energy spacetime description follows a Jacobson-type Einstein branch, why the gauge structure is
\[
\frac{\SU(3)\times\SU(2)\times\U(1)}{\mathbb Z_6},
\]
why there are three colors and three generations, why hypercharges take the observed rational values, why the electroweak and flavor scales are what they are, and why dark-sector phenomenology and the cosmological constant enter in the particular form observed.

This paper presents a single framework that addresses this list as a conditional reconstruction program. The framework is Observer-Patch Holography (OPH). Its starting point is an observer-centric organization of physical data: each observer has access only to a patch algebra, neighboring observers must agree on overlaps, and global physical law is recovered from the resulting consistency conditions. The same formal structure that produces schedule-independent overlap consistency then feeds a Lorentzian modular-phase classification, a Jacobson-type Einstein branch conditional on fixed-cap stationarity and the null-stress premise, compact gauge reconstruction in a bosonic internal-gauge sector, charge quantization, and the low-energy architecture selected by the extended MAR admissibility package.

Two decisions shape the present version.

\begin{enumerate}[leftmargin=2em]
\item The compact note presents the main theorem statements in one place, but some quantitative conventions and implementation details live in repository supplements.\footnote{Questions about ultimate purpose are deferred until the corresponding physical question has been posed with enough precision to deserve an answer.}
\item The paper separates exact structural theorems, conditional scaling-limit branches, calibration-sector outputs, and phenomenological continuations. This keeps strong claims visible and weak claims identifiable.
\end{enumerate}

The framework is intended as a candidate fundamental theory in the precise sense relevant for physics: general relativity and the Standard Model appear as effective sectors derived from a smaller axiom set plus an explicit refinement/scaling-limit package. The paper does not use simulation language, observer continuation stories, or other non-falsifiable additions.

\subsection*{Prioritized achievements}

The results emphasized in this paper are ordered below by a combination of impact and defensibility.

\begin{enumerate}[leftmargin=2em]
\item A conditional Jacobson-type Einstein branch is recovered from observer-overlap modular geometry, an explicit fixed-cap generalized-entropy stationarity premise, and an explicit null-modular / null-stress scaling-limit package.
\item The exact Standard Model gauge structure
\[
\frac{\SU(3)\times\SU(2)\times\U(1)}{\mathbb Z_6}
\]
is recovered in the bosonic compact-gauge branch once the extended MAR admissibility package is imposed.
\item On the realized one-generation chiral matter plus one-Higgs package, anomaly constraints and Yukawa invariance fix the hypercharge lattice, and MAR then selects three colors and three generations on the realized admissible branch.
\item The current quantitative implementation uses two external continuous configuration inputs, $P$ and $N_{\mathrm{scr}}$.
\item The Higgs/top critical-surface branch adds no extra continuous fit once the gauge trajectory and scale-setting branch are fixed, while the charged-lepton/Koide continuation uses extra phenomenological structure beyond the core theorem package.
\item A separate input-dependent capacity corollary ties the cosmological-constant branch to global screen capacity rather than local vacuum energy.
\item Absence of gauge-mediated proton decay follows from product-group structure.
\item Dark-sector, baryogenesis, string, and spectroscopy branches are kept visible, but only as program-level or conjectural phenomenology.
\item Objective physical law is identified with the unique normal form of the overlap-consistency problem.
\end{enumerate}

\section{Results at a Glance}

This section is the authoritative status ledger for the compact note and the sync rule for mirrored summary surfaces. The three-phase boundary used throughout is
\[
\text{Phase I}=(D1\text{--}D5)\cup(D7\text{--}D9),\qquad
\text{Phase II}=D6\cup D10\cup D11,\qquad
\text{Phase III}=D12.
\]
Phase I is the recovered-core claim set. Phase II contains the separate input-dependent capacity corollary and the current supplement-backed numerical implementation. Phase III collects deferred continuations. The canonical five-axiom basis, the two external continuous inputs, and the technical-premise labels used below are fixed in Section~\ref{sec:premise-ledger}. Each row names a node in the reconstruction program, records its immediate OPH-side parents, separates imported standard mathematics from extra non-axiom physical premises or external inputs, and states the resulting claim tier.

\begingroup
\scriptsize
\setlength{\tabcolsep}{2pt}
\begin{longtable}{@{}>{\raggedright\arraybackslash}p{0.05\textwidth}>{\raggedright\arraybackslash}p{0.16\textwidth}>{\raggedright\arraybackslash}p{0.14\textwidth}>{\raggedright\arraybackslash}p{0.13\textwidth}>{\raggedright\arraybackslash}p{0.20\textwidth}>{\raggedright\arraybackslash}p{0.16\textwidth}@{}}
\caption{Authoritative dependency checklist for the compact OPH program.}\\
\toprule
Node & Output & Immediate OPH ingredients & Standard mathematics used & Extra non-axiom physical premises / external inputs & Claim tier \\
\midrule
\endhead
\bottomrule
\endfoot
\textbf{D1} & Unique schedule-independent normal form (Theorem~\ref{thm:confluence}) & overlap-consistency problem on a finite patch net & Newman's lemma for terminating locally confluent rewrite systems & finite state spaces, termination, and local confluence & Phase I structural theorem \\
\textbf{D2} & Collar Markov/entropy split and \(G\)-dictionary setup (Theorem~\ref{thm:markovcollar}, Propositions~\ref{prop:approxrecover}, \ref{prop:splice}, \ref{prop:gensplit}, Definition~\ref{def:newton}) & Axioms~\ref{ax:screen}--\ref{ax:entropy} & HJPW Markov structure theorem; Petz and Fawzi--Renner recovery theory & exact block-factorized identities require exact Markovity or an explicitly stated idealized limit; the \(A/(4G)\) step is the coarse-grained dictionary of Definition~\ref{def:newton} & Phase I structural lemma/proposition layer \\
\textbf{D3} & Lorentz branch \(\Conf^+(S^2)\cong \mathrm{SO}^+(3,1)\) (Theorem~\ref{thm:bw}, Corollary~\ref{cor:lorentz}) & Axioms~\ref{ax:screen}--\ref{ax:entropy} & Bisognano--Wichmann modular geometry; conformal-group classification & T4, T5, and T8 of Assumption~\ref{ass:technical} & Phase I conditional scaling-limit theorem \\
\textbf{D4} & Null modular bridge to \(T_{kk}\) (Propositions~\ref{prop:nullec}, \ref{prop:n3}; Corollaries~\ref{cor:n1}, \ref{cor:n2}; Theorem~\ref{lem:density}, Lemma~\ref{lem:translation}; Theorem~\ref{thm:stress}) & D2+D3 and Axioms~\ref{ax:screen}--\ref{ax:entropy} & Borchers--Wiesbrock half-sided modular inclusion; distributional differentiation on half-lines & T1--T3, T6, T8, and the relativistic null-stress identification premise & Phase I conditional bridge theorem \\
\textbf{D5} & Jacobson-type Einstein branch: rest-frame relation plus conditional tensor upgrade (Lemma~\ref{lem:smallball}, Theorem~\ref{thm:einstein}, Corollary~\ref{cor:einstein}) & D3+D4 and Axioms~\ref{ax:screen}--\ref{ax:entropy} & Jacobson small-ball first law; null-to-tensor reconstruction modulo a metric term & T9 fixed-cap stationarity; locally Lorentzian \(d=4\) scaling regime; small-ball constancy assumptions; all local directions/reference states for the tensor upgrade & Phase I conditional scaling-limit theorem/corollary \\
\textbf{D6} & Capacity branch for \(\Lambda\) and the separate order-of-magnitude neutrino estimate (Proposition~\ref{prop:lambda}, Corollary~\ref{cor:lambda}) & D5 leaves a \(+\Lambda g_{ab}\) ambiguity & de Sitter entropy relation and dimensional analysis & external input \(N_{\mathrm{scr}}\) and the screen-capacity identification & Phase II input-dependent corollary / estimate \\
\textbf{D7} & Compact gauge reconstruction (Theorem~\ref{thm:compactgauge}) & Axioms~\ref{ax:screen}--\ref{ax:entropy} & Doplicher--Roberts / Tannaka reconstruction & T7 and T10--T12; rigid symmetric bosonic refinement-stable sector colimit with fiber functor & Phase I conditional structural theorem \\
\textbf{D8} & Product gauge structure up to finite quotient (Lemmas~\ref{lem:mincontent}--\ref{lem:commutant}, Theorem~\ref{thm:sm}) & D7 + Axiom~\ref{ax:mar} & compact Lie representation classification; Schur's lemma & connected positive-dimensional Lie admissible class; one connected abelian factor; faithful action on the minimal coupled carrier; vanishing obstruction \([z]=0\) & Phase I realized-branch theorem \\
\textbf{D9} & Realized Standard Model quotient, hypercharges, \(N_c=3\), \(N_g=3\), and product-group corollaries (Theorem~\ref{thm:hypercharge}, Corollaries~\ref{cor:ng}, \ref{cor:nc}, Proposition~\ref{prop:z6}) & D8 & anomaly-cancellation algebra; Witten global-anomaly argument; CKM CP counting & realized one-generation chiral matter plus one Higgs package; asymptotic-freedom and MAR admissibility premises & Phase I realized-branch theorem/corollary chain \\
\textbf{D10} & Gauge/electroweak numerical closure of the current supplement-backed implementation & D9 + external input \(P\) & RG evolution and matching conventions supplied in the supplement & pixel constraint; edge heat-kernel closure; supplement-backed beta-function and threshold conventions & Phase II calibration sector \\
\textbf{D11} & Higgs/top critical-surface branch of the current supplement-backed implementation & D10 & RG transport and pole-mass conversion supplied in the supplement & critical-surface condition \(\lambda=\beta_\lambda=0\) at \(M_U\); no new continuous fit once D10 is fixed & Phase II supplement-backed quantitative branch \\
\textbf{D12} & Charged-lepton/Koide, texture, dark-sector, baryogenesis, black-hole spectroscopy, and string/worldsheet continuations & various subsets of D6, D9, D10, and D11 & branch-specific EFT and phenomenological manipulations & additional ans\"atze such as the uniform \(\mathbb Z_6\) center-label ensemble, discrete texture choices, dark-sector response assumptions, or discrete-horizon assumptions & Phase III phenomenological continuation / program-level branch \\
\end{longtable}
\endgroup

Accordingly, the recovered-core structural program of the compact note is D1--D5 together with D7--D9. D6 is the separate Phase-II input-dependent corollary; D10--D11 are the current Phase-II supplement-backed numerical implementation; and D12 is the Phase-III continuation layer. The only external continuous inputs are \(N_{\mathrm{scr}}\) in D6 and \(P\) in D10, together with descendants built from those branches.
\section{Five Axioms, Two External Continuous Inputs in the Current Quantitative Implementation, and Technical Premises}\label{sec:premise-ledger}

This section is the authoritative premise ledger and dependency map for the compact note. The OPH basis used below is exactly the five core axioms, the two external continuous inputs of the current quantitative implementation, and the theorem-local technical premises collected afterward. Older A1--A4 / B / MX / LR / R0--R1 packaging is retained elsewhere only as shorthand for pieces of this canonical package, not as a competing axiom basis.

The formulation uses two external continuous configuration inputs in the current quantitative implementation:
\begin{align}
P &\equiv a_{\mathrm{cell}}/\ell_P^2 = 1.63094,\\
N_{\mathrm{scr}} &\equiv \log \dim \mathcal H_{\mathrm{tot}} \sim 10^{122}.
\end{align}
In the current implementation, \(P\) is calibrated from the gauge sector and \(N_{\mathrm{scr}}\) is inferred from the observed cosmological constant. The first sets the particle-physics scale. The second enters the cosmological-capacity branch and a separate order-of-magnitude neutrino estimate.

\begin{axiom}[Screen Net]\label{ax:screen}
Physical data is organized on a horizon screen \(S^2\) carrying a net of local algebras
\[
P \longmapsto \mathcal A(P)
\]
for connected patches \(P\subset S^2\), with isotony
\[
P\subset Q \implies \mathcal A(P)\subset \mathcal A(Q).
\]
\end{axiom}

\begin{axiom}[Overlap Consistency]\label{ax:overlap}
For overlapping patches \(P_1\cap P_2\neq\varnothing\), the local states induced on the shared algebra agree:
\[
\omega_{P_1}|_{\mathcal A(P_1\cap P_2)}
=
\omega_{P_2}|_{\mathcal A(P_1\cap P_2)}.
\]
\end{axiom}

\begin{axiom}[Local MaxEnt and Refinement Stability]\label{ax:maxent}
At the UV scale the realized state maximizes entropy subject to a finite family of gauge-invariant local constraints. The resulting family of states is stable under coarse-graining, so symmetry-allowed relevant operators are not held at zero by unexplained fine tuning.
\end{axiom}

\begin{axiom}[Recoverable Generalized Entropy]\label{ax:entropy}
A generalized entropy functional exists on caps,
\[
S_{\mathrm{gen}}(C)=S_{\mathrm{bulk}}(C)+\langle L_C\rangle,
\]
where \(L_C\) is a positive edge-center entropy functional. In the semiclassical scaling branch, its leading coarse-grained contribution is identified with \(A(\partial C)/(4G)\). The functional obeys quantum focusing on null generators, and collar tripartitions have small conditional mutual information with controlled recovery maps \cite{LiebRuskai,Petz86,Petz88,FawziRenner}.
\end{axiom}

\begin{axiom}[Minimal Admissible Realization]\label{ax:mar}
Among realized sector packages \(\mathfrak S\) consisting of the connected Lie gauge-sector image relevant in the low-energy EFT, its admissible light chiral matter content, and one Higgs doublet, and which are loop-coherent, anomaly-free, refinement-stable with light chiral matter, single-Higgs Yukawa-completable with one connected abelian charge factor acting nontrivially on the coupled carrier, intrinsically CP-capable, and weak-sector UV-completable, the realized package is the lexicographically minimal one under
\[
C(\mathfrak S)=\bigl(\chi_{\mathrm{cpl}},\,N_{\mathrm{nonab}},\,N_c,\,N_g\bigr).
\]
\(\mathfrak S\) is the sector package on which MAR acts; it is not the bare tensor category alone. MAR is therefore an explicit structural-economy axiom on already-admissible realized low-energy branches, not a theorem derived from the preceding axioms. \(\chi_{\mathrm{cpl}}\) is the coupled carrier dimension: the smallest unitary carrier containing a common irreducible block on which the admissible pseudoreal and complex nonabelian charge types both act nontrivially. This is stronger than the abstract minimal faithful representation dimension.
\end{axiom}

\begin{assumption}[Technical Premises]\label{ass:technical}
The continuum reductions additionally use:
\begin{enumerate}[label=\textbf{T\arabic*},leftmargin=2.3em]
\item type-I local algebras at the regulator scale and boundary gauge invariance on collars;
\item finite-velocity quasi-local propagation at the UV scale;
\item quantitative collar-mixing / endpoint-Lipschitz control for null-interval modular Hamiltonians on bounded intervals;
\item cap isotropy / geometric modular action around the entangling circle in the refinement-scaling regime;
\item Euclidean modular regularity, fixing the angular period to \(2\pi\);
\item a standard half-sided modular inclusion premise for the null half-line blow-up algebras;
\item vanishing gluing obstruction \([z]=0\) whenever global sector transportability is invoked;
\item all Lorentz/null-modular/Einstein statements are claims about a controlled refinement/scaling limit of the regulator net, not literal exactness at fixed finite cutoff;
\item fixed-cap generalized-entropy stationarity for the admissible first-variation class used in the Jacobson branch;
\item symmetric braiding in the \(3+1\)D EFT branch;
\item a bosonic Tannakian fiber-functor premise, or else an explicit super-Tannakian fork;
\item whenever compact gauge reconstruction is invoked, the sector category means the refinement-stable directed colimit of transportable edge sectors with objectwise finite-dimensional fibers, not the block list of one fixed cutoff algebra.
\end{enumerate}
\end{assumption}

\begin{remark}[Legacy shorthand]
Historical labels are read as follows in the compact note and in mirrored surfaces: A1/A2 mean Screen Net and Overlap Consistency; B together with refinement stability means Axiom~\ref{ax:maxent}; A3 means Axiom~\ref{ax:entropy}; A4 together with the old MX / collar terminology means the recoverability and collar-control ingredients explicitly cited from Axiom~\ref{ax:entropy} and Assumption~\ref{ass:technical}; and R0/R1, LR, and \( [z]=0 \) are theorem-local technical premises rather than additional axioms.
\end{remark}

\begin{remark}
The five axioms fix the OPH core, but the compact note's conditional reconstruction program also uses the scaling-limit, categorical, and stationarity premises collected in Assumption~\ref{ass:technical}. In particular, the Lorentz/null-modular/Einstein and compact-gauge statements below are scaling-limit claims, not fixed-cutoff matrix-algebra identities. These premises are not derived in the compact note; they are target regularity properties of the intended Lorentzian EFT phase and candidate UV completions.
\end{remark}

\begin{remark}[Log convention]
Unless an explicit base is written, logarithms and entropies in this compact note use natural logs (nats). Base-2 logarithms are written as \(\log_2\) and quoted in bits.
\end{remark}

\subsection*{Authoritative dependency DAG}

The table below is the authoritative dependency map for the compact note. The ``Immediate OPH ingredients'' column records only parents internal to the OPH program. Imported standard mathematics and extra non-axiom physical premises or external inputs are stated separately rather than hidden inside the word ``derived.''

\begingroup
\scriptsize
\setlength{\tabcolsep}{2pt}
\begin{longtable}{@{}>{\raggedright\arraybackslash}p{0.06\textwidth}>{\raggedright\arraybackslash}p{0.20\textwidth}>{\raggedright\arraybackslash}p{0.14\textwidth}>{\raggedright\arraybackslash}p{0.13\textwidth}>{\raggedright\arraybackslash}p{0.21\textwidth}>{\raggedright\arraybackslash}p{0.12\textwidth}@{}}
\toprule
Node & Compact theorem labels and output & Immediate OPH ingredients & Standard mathematics used & Extra non-axiom physical premises / external inputs & Status \\
\midrule
\endhead
\bottomrule
\endfoot
\textbf{D1} & Theorem~\ref{thm:confluence}: unique normal form and schedule independence & overlap-consistency problem on a finite patch net & Newman's lemma for terminating locally confluent rewrite systems & finite state spaces, termination, and local confluence & structural theorem \\
\textbf{D2} & Theorem~\ref{thm:markovcollar}, Propositions~\ref{prop:approxrecover}, \ref{prop:splice}, \ref{prop:gensplit}, Definition~\ref{def:newton}: collar Markov/entropy layer & Axioms~\ref{ax:screen}--\ref{ax:entropy} & HJPW Markov structure theorem; Petz and Fawzi--Renner recovery theory & exact Markov identities require exact Markovity or an explicitly stated idealized recoverability limit; the \(A/(4G)\) identification is the coarse-grained dictionary of Definition~\ref{def:newton} & structural lemma/proposition layer \\
\textbf{D3} & Theorem~\ref{thm:bw}, Corollary~\ref{cor:lorentz}: geometric modular flow and Lorentz branch & Axioms~\ref{ax:screen}--\ref{ax:entropy} & Bisognano--Wichmann modular geometry; \(\Conf^+(S^2)\cong \mathrm{SO}^+(3,1)\) & T4, T5, and T8 of Assumption~\ref{ass:technical} & conditional scaling-limit theorem \\
\textbf{D4} & Propositions~\ref{prop:nullec}, \ref{prop:n3}; Corollaries~\ref{cor:n1}, \ref{cor:n2}; Theorem~\ref{lem:density}, Lemma~\ref{lem:translation}; Theorem~\ref{thm:stress}: null modular bridge to \(T_{kk}\) & D2+D3 and Axioms~\ref{ax:screen}--\ref{ax:entropy} & Borchers--Wiesbrock half-sided modular inclusion; distributional differentiation on half-lines & T1--T3, T6, T8, and the relativistic null-stress identification premise & conditional bridge theorem \\
\textbf{D5} & Lemma~\ref{lem:smallball}, Theorem~\ref{thm:einstein}, Corollary~\ref{cor:einstein}: Jacobson-type Einstein branch via a rest-frame relation plus conditional tensor upgrade & D3+D4 and Axioms~\ref{ax:screen}--\ref{ax:entropy} & Jacobson small-ball first law; Lemma~\ref{lem:nullreconstruct} for the null-to-tensor upgrade modulo a metric term & T9 fixed-cap stationarity; locally Lorentzian \(d=4\) scaling regime; small-ball constancy assumptions; all local directions/reference states for the tensor upgrade & conditional scaling-limit theorem/corollary \\
\textbf{D6} & Proposition~\ref{prop:lambda}, Corollary~\ref{cor:lambda}: capacity branch for \(\Lambda\), plus the separate neutrino estimate & D5, which already leaves the \(+\Lambda g_{ab}\) ambiguity & de Sitter entropy relation and dimensional analysis & external input \(N_{\mathrm{scr}}\) and the screen-capacity identification & input-dependent corollary / estimate \\
\textbf{D7} & Theorem~\ref{thm:compactgauge}: compact gauge reconstruction & Axioms~\ref{ax:screen}--\ref{ax:entropy} & Doplicher--Roberts / Tannaka reconstruction & T7 and T10--T12; rigid symmetric bosonic refinement-stable sector colimit with fiber functor & conditional structural theorem \\
\textbf{D8} & Lemmas~\ref{lem:mincontent}--\ref{lem:commutant}, Theorem~\ref{thm:sm}: product gauge structure up to finite quotient & D7 + Axiom~\ref{ax:mar} & compact Lie representation classification; Schur's lemma & connected positive-dimensional Lie admissible class; one connected abelian factor; faithful action on the minimal coupled carrier; vanishing obstruction \( [z]=0 \) & realized-branch theorem \\
\textbf{D9} & Theorem~\ref{thm:hypercharge}, Corollaries~\ref{cor:ng}, \ref{cor:nc}, Proposition~\ref{prop:z6}: realized Standard Model quotient, hypercharges, \(N_c=3\), \(N_g=3\), and product-group corollaries & D8 & anomaly-cancellation algebra; Witten global-anomaly argument; CKM CP counting & realized one-generation chiral matter plus one Higgs package; asymptotic-freedom and MAR admissibility premises & realized-branch theorem/corollary chain \\
\textbf{D10} & quantitative gauge/electroweak closure of Section~7 & D9 + external input \(P\) & RG evolution and matching conventions supplied in the supplement & pixel constraint; edge heat-kernel closure; supplement-backed beta-function and threshold conventions & calibration sector \\
\textbf{D11} & Higgs/top subsection of Section~7 & D10 & RG transport and pole-mass conversion supplied in the supplement & critical-surface condition \(\lambda=\beta_\lambda=0\) at \(M_U\); no new continuous fit once D10 is fixed & supplement-backed quantitative branch \\
\textbf{D12} & Sections~8--10 program continuations & various subsets of D6, D9, D10, and D11 & branch-specific EFT and phenomenological manipulations & additional ans\"atze such as the uniform \(\mathbb Z_6\) center-label ensemble, texture choices, dark-sector response assumptions, or discrete-horizon assumptions & phenomenological continuation / program-level branch \\
\end{longtable}
\endgroup

No theorem in the compact note is meant to hide standard mathematical imports or extra physical premises inside the phrase ``from the axioms.'' The authoritative unpacking of every such claim is the dependency DAG above. The only rows that use the two external continuous inputs are D6, D10, the separate order-of-magnitude neutrino estimate attached to D6, and downstream branches built from them.

\begin{theorem}[Summary theorem for the recovered relativity-plus-Standard-Model core]\label{thm:master}
This theorem packages the recovered-core nodes D1--D5 and D7--D9 of the authoritative dependency DAG above. Node D6 is a separate input-dependent corollary, while D10--D12 remain outside this theorem's scope. Assume Axioms~\ref{ax:screen}--\ref{ax:entropy}, Axiom~\ref{ax:mar}, Assumption~\ref{ass:technical}, and the hypotheses of Theorems~\ref{thm:confluence}, \ref{thm:bw}, \ref{thm:stress}, \ref{thm:einstein}, \ref{thm:compactgauge}, \ref{thm:sm}, and \ref{thm:hypercharge}, together with Corollaries~\ref{cor:ng}, \ref{cor:nc}, and Proposition~\ref{prop:z6}. Then, in the controlled refinement/scaling regime singled out by Assumption~\ref{ass:technical}:
\begin{enumerate}[label=(\roman*),leftmargin=2em]
\item overlap repair admits a unique schedule-independent normal form;
\item cap modular flow is geometric in the assumed modular phase and yields the connected Lorentz group
\[
\Conf^+(S^2)\cong \mathrm{SO}^+(3,1);
\]
\item admissible fixed-cap stationarity together with the null modular bridge and null-stress premise yield the Jacobson-type rest-frame relation
\[
\delta\!\left(G_{00}+\Lambda g_{00}\right)=8\pi G\,\delta\langle T_{00}\rangle
\]
at the cap center in the diamond rest frame; if that rest-frame relation holds for all local directions and reference states in the scaling regime, then the null-to-tensor upgrade gives
\[
G_{ab}+\Lambda g_{ab}=8\pi G\,\langle T_{ab}\rangle;
\]
\item the refinement-stable transportable edge-sector category reconstructs a compact gauge group, and the realized connected gauge structure has the form
\[
\frac{\SU(3)\times\SU(2)\times\U(1)}{\Gamma}
\]
for some finite central subgroup \(\Gamma\).
\end{enumerate}
After the hypercharge lattice, generation-count, color-count, and trivial-action quotient steps, the finite quotient is fixed to \(\Gamma=\mathbb Z_6\), so the realized gauge structure on the realized MAR-admissible branch is
\[
\frac{\SU(3)\times\SU(2)\times\U(1)}{\mathbb Z_6},
\qquad
N_c=3,
\qquad
N_g=3,
\]
with the exact Standard Model hypercharge lattice on the realized matter package.
\end{theorem}

\begin{proof}
Items (i)--(iii) are collected from Theorem~\ref{thm:confluence}, Corollary~\ref{cor:lorentz}, Theorem~\ref{thm:einstein}, and Corollary~\ref{cor:einstein}. The compact gauge reconstruction is Theorem~\ref{thm:compactgauge}; the product gauge structure up to finite quotient is Theorem~\ref{thm:sm}; and the final identification of the exact Standard Model quotient is Theorem~\ref{thm:hypercharge} together with Corollaries~\ref{cor:ng}, \ref{cor:nc}, and Proposition~\ref{prop:z6}.
\end{proof}

Theorem~\ref{thm:master} therefore packages only the recovered relativity-plus-Standard-Model core: D1--D5 together with D7--D9. The D6 capacity relation is a separate input-dependent corollary; the quantitative calibration sector, Higgs/top branch, and phenomenological continuations remain outside its scope.

\subsection*{Usually Postulated Elsewhere, Recovered or Continued Here with Explicit Status}

\begingroup
\footnotesize
\begin{tabular}{@{}>{\raggedright\arraybackslash}p{0.28\textwidth}>{\raggedright\arraybackslash}p{0.26\textwidth}>{\raggedright\arraybackslash}p{0.36\textwidth}@{}}
\toprule
Structure & Common treatment & OPH treatment \\
\midrule
Lorentz kinematics & background symmetry or starting axiom & conditional scaling-limit branch from geometric modular flow on screen caps \\
Einstein dynamics & fundamental field equation & conditional first-variation relation with tensor upgrade from generalized entropy, null modular data, and entanglement equilibrium \\
Gauge group & model input & conditional compact group reconstruction from edge sectors; exact SM quotient selected on the realized MAR-admissible branch \\
Hypercharge lattice & matter-assignment input & solved from anomaly cancellation and Yukawa invariance on the realized one-generation chiral matter plus one-Higgs package \\
Color and generation count & empirical input & fixed on the realized MAR-admissible branch by anomaly, CP capability, asymptotic freedom, and minimality \\
Flavor hierarchy unit & model-dependent small parameter & phenomenological ansatz \(\varepsilon=1/6\) motivated by a uniform \(\mathbb Z_6\) center-label ensemble \\
Koide phase & phenomenological fit & continuation ansatz tied to the previously fixed discrete data as \(\delta=2/9\) \\
Cosmological constant & local vacuum-energy puzzle & global screen-capacity term after the input-dependent capacity identification \\
\bottomrule
\end{tabular}
\endgroup
\section{Observer-Overlap Consistency as a Fixed-Point Problem}

The discrete overlap problem is the mathematical skeleton of the framework.

\begin{definition}[Patch Net]\label{def:patchnet}
Let $G=(V,E)$ be a finite connected graph. Each vertex $i\in V$ carries a finite local state space $S_i$. For each edge $e=\{i,j\}$ let $I_e$ be an interface alphabet and let
\[
\pi_{i,e}:S_i\to I_e,
\qquad
\pi_{j,e}:S_j\to I_e
\]
be interface projections. The global state space is
\[
\Sigma=\prod_{i\in V} S_i,
\]
and the consistency set is
\[
C=\bigl\{s\in\Sigma:\pi_{i,e}(s_i)=\pi_{j,e}(s_j)\ \text{for all }e=\{i,j\}\in E\bigr\}.
\]
\end{definition}

This finiteness is a regulator-level assumption for the discrete patch-net normal-form theorem; it is not the continuum limit statement.

Define the inconsistency potential
\[
\Phi(s)=\sum_{e=\{i,j\}\in E}w_e\, d_e\!\bigl(\pi_{i,e}(s_i),\pi_{j,e}(s_j)\bigr),
\]
with $w_e>0$ and $d_e(a,b)=0$ iff $a=b$.

\begin{definition}[Local Repair Law]
A law $\lambda$ is a family of local repair maps
\[
T_i^\lambda:\Sigma\to\Sigma
\]
changing only patch $i$ or a bounded neighborhood of $i$.
\end{definition}

\begin{theorem}[Asynchronous Fixed-Point Law]\label{thm:confluence}
Assume:
\begin{enumerate}[label=(\roman*),leftmargin=2em]
\item normal forms of the repair relation are exactly the globally consistent states $C$;
\item the repair relation terminates;
\item the repair relation is locally confluent.
\end{enumerate}
Then every $s\in\Sigma$ has a unique normal form
\[
\nf_\lambda(s)\in C,
\]
and every maximal asynchronous repair sequence terminates at that same state.
A sufficient condition for (ii) is that each $S_i$ is finite and every enabled repair strictly decreases $\Phi$.
\end{theorem}

\begin{proof}
Termination is assumed. Local confluence plus termination imply confluence by Newman's lemma. The resulting normal form is unique. By assumption it is a globally consistent state.
\end{proof}

\begin{corollary}[Schedule Independence]
If observables are evaluated on $\nf_\lambda(s)$, the resulting physical law is independent of update order.
\end{corollary}

\begin{theorem}[Cycle Obstruction / Holonomy Criterion]\label{thm:holonomy}
Let $A$ be an abelian group. For each oriented edge $e:u\to v$ assign $b_e\in A$ and consider the affine overlap equations
\[
x_v-x_u=b_e.
\]
A global solution exists if and only if the signed sum of $b_e$ vanishes on every cycle.
\end{theorem}

\begin{proof}
Summing the edge equations around a cycle gives necessity. For sufficiency, fix a root and define each $x_v$ by summing labels along a path from the root to $v$. Vanishing cycle sums make this independent of the chosen path.
\end{proof}

\begin{proposition}[Gauge Quotient]\label{prop:gauge}
Suppose a local gauge group $\Gamma=\prod_i \Gamma_i$ acts on $\Sigma$ while leaving all interface data invariant, and the repair maps satisfy
\[
T_i^\lambda(\gamma\cdot s)=\gamma\cdot T_i^\lambda(s).
\]
Then
\[
\nf_\lambda(\gamma\cdot s)=\gamma\cdot \nf_\lambda(s).
\]
Hence gauge-invariant observables are unique on the quotient.
\end{proposition}

\begin{proof}
Gauge covariance maps repair sequences to repair sequences. Since the normal form is unique by Theorem~\ref{thm:confluence}, the normal form of $\gamma\cdot s$ must be $\gamma\cdot \nf_\lambda(s)$.
\end{proof}

These three results already encode much of the eventual physical interpretation. Objectivity becomes confluence, gauge symmetry becomes quotient invariance, and stable defects are represented by nontrivial overlap holonomy classes.

\section{Collars, Edge Centers, and Generalized Entropy}

The fixed-point picture explains why overlap consistency produces gauge quotients, but older presentations still inserted a preferred boundary gluing group when the collar-center decomposition was proved. The point of this section is to remove that extra model dependence at fixed regulator scale. Once one chooses any finite regulator cover of a cap neighborhood, overlap-consistent rechartings of the cut data form a compact unitary transition system. On the ordinary or central-defect branch, that transition system reduces to a compact boundary group or compact central extension, and the collar-center block decomposition follows from its invariant theory. The only fixed-cutoff obstruction to that ordinary-group statement is a genuinely noncentral 2-group defect, which we keep explicit rather than bury in packaging.

\begin{proposition}[Derived regulator gluing datum]\label{prop:regulatedrealization}
Fix a regulator-scale finite patch cover of a cap neighborhood and let \(R\) be any finite union of regulator cells. Assume only the finite patch-net presentation implicit in Definition~\ref{def:patchnet} together with overlap consistency on common interfaces. Then:
\begin{enumerate}[leftmargin=2em]
\item each regulator cell \(i\) with finite local state space \(S_i\) can be Hilbertized as \(\tilde{\mathcal H}_i\cong \mathbb C^{|S_i|}\), so the extended algebra before overlap quotienting is the finite type-I algebra
\[
\widetilde{\mathcal A}(R)=\mathcal B(\tilde{\mathcal H}_R),
\qquad
\tilde{\mathcal H}_R:=\bigotimes_{i\subset R}\tilde{\mathcal H}_i;
\]
\item for each connected boundary component \(\Sigma\subset\partial R\), after choosing a reference cut presentation \(\tilde{\mathcal H}_\Sigma\), every overlap-consistent change of local chart along \(\Sigma\) is implemented by a unitary on \(\tilde{\mathcal H}_\Sigma\), and the compact closure of the subgroup generated by all such recharting unitaries is a compact boundary redundancy group
\[
K_\Sigma\subset U(\tilde{\mathcal H}_\Sigma);
\]
before the choice of reference chart, the same data form a compact unitary groupoid of overlap-preserving transitions;
\item if triple-overlap defects are central, the projective composition law of item 2 lifts to an honest action of a compact central extension \(\widehat K_\Sigma\); a genuinely noncentral defect is the only obstruction to reducing the transition system to an ordinary compact group action;
\item in the invariant-state realization of the quotient, physical observables are the chart-independent endomorphisms, equivalently the fixed-point algebra
\[
\mathcal A(R)=\widetilde{\mathcal A}(R)^{\widehat K_{\partial R}},
\qquad
\widehat K_{\partial R}:=\prod_{\Sigma\subset\partial R}\widehat K_\Sigma,
\]
\noindent with the convention \(\widehat K_\Sigma=K_\Sigma\) whenever no central extension is needed.
\end{enumerate}
Historical labels \(R0\) and \(R1\) are shorthand for items 1 and 4 on that branch; they are not extra axioms beyond the finite regulator presentation plus overlap redundancy.
\end{proposition}

\begin{proof}[Proof sketch]
Item 1 is the Hilbertization of finite sets. For item 2, each overlap-consistent recharting acts on a finite-dimensional matrix algebra, and every \(^*\)-automorphism of such an algebra is inner. After transporting all local overlap presentations to a chosen reference chart, each recharting is therefore implemented by a unitary on the cut Hilbert space. The subgroup generated by those unitaries has compact closure inside a finite-dimensional unitary group, which yields \(K_\Sigma\); without fixing the reference chart one has the corresponding compact unitary groupoid. Item 3 is the usual projective-versus-central-extension lift for a central 2-cocycle. Item 4 is the quotient statement: physical observables are precisely the endomorphisms insensitive to recharting of the boundary data, hence the fixed points of the derived boundary action. If the defect is genuinely noncentral, the quotient is still well defined, but the correct bookkeeping object is the crossed-module or higher-gauge data rather than an ordinary compact group.
\end{proof}

\begin{theorem}[Derived edge-center collar decomposition and exact Markov normal form]\label{thm:markovcollar}
Let \(A\)-\(B\)-\(D\) be a collar tripartition realized at fixed regulator scale on the ordinary or central-defect branch of Proposition~\ref{prop:regulatedrealization}. Write \(B=B_L\cup B_R\) with common interface \(\Sigma=\partial C\), and let \(\widehat K_\Sigma\) be the derived compact boundary action attached to that cut, with \(\widehat K_\Sigma=K_\Sigma\) on the ordinary branch. In the invariant-state realization,
\[
\mathcal H_B=(\tilde{\mathcal H}_{B_L}\otimes \tilde{\mathcal H}_{B_R})^{\widehat K_\Sigma}.
\]
Decompose the left boundary data as
\[
\tilde{\mathcal H}_{B_L}\cong \bigoplus_\alpha W_\alpha\otimes \mathcal H_{b_L^\alpha},
\]
with \(W_\alpha\) irreducible \(\widehat K_\Sigma\)-modules. Because the right half-collar carries the inverse overlap transport across the same cut, it decomposes contragrediently:
\[
\tilde{\mathcal H}_{B_R}\cong \bigoplus_\beta W_\beta^*\otimes \mathcal H_{b_R^\beta}.
\]
Then
\[
\mathcal H_B\cong\bigoplus_\alpha \mathcal H_{b_L^\alpha}\otimes \mathcal H_{b_R^\alpha},
\qquad
Z(\mathcal A(B))=\bigoplus_\alpha \mathbb C\,\mathbf 1_\alpha.
\]
If, in addition, the reference state on \(A\cup B\cup D\) is exact Markov,
\[
I(A:D|B)_\rho=0,
\]
or one passes to an explicitly stated idealized recoverability limit that reduces to it, then the state decomposes as
\[
\rho_{ABD}
=
\bigoplus_\alpha p_\alpha\,
\rho^{(\alpha)}_{A b_L^\alpha}\otimes \rho^{(\alpha)}_{b_R^\alpha D}.
\]
Thus the collar-center block decomposition is forced by overlap consistency plus the derived regulator package on the ordinary or central-defect branch; exact Markov factorization is an additional state condition, not a consequence of the block decomposition alone.
\end{theorem}

\begin{proof}
By Proposition~\ref{prop:regulatedrealization}, one may realize the physical collar states as the \(\widehat K_\Sigma\)-invariant subspace of \(\tilde{\mathcal H}_{B_L}\otimes \tilde{\mathcal H}_{B_R}\). Complete reducibility of finite-dimensional unitary representations gives the displayed decomposition of \(\tilde{\mathcal H}_{B_L}\), and the right half-collar carries the inverse transport law across the same cut, hence the contragredient decomposition of \(\tilde{\mathcal H}_{B_R}\). Therefore
\[
\tilde{\mathcal H}_{B_L}\otimes \tilde{\mathcal H}_{B_R}
\cong
\bigoplus_{\alpha,\beta}
(W_\alpha\otimes W_\beta^*)\otimes
(\mathcal H_{b_L^\alpha}\otimes \mathcal H_{b_R^\beta}).
\]
Taking \(\widehat K_\Sigma\)-invariants gives
\[
\mathcal H_B
\cong
\bigoplus_{\alpha,\beta}
(W_\alpha\otimes W_\beta^*)^{\widehat K_\Sigma}\otimes
(\mathcal H_{b_L^\alpha}\otimes \mathcal H_{b_R^\beta}).
\]
By Schur's lemma,
\[
(W_\alpha\otimes W_\beta^*)^{\widehat K_\Sigma}
\cong
\begin{cases}
\mathbb C, & \alpha=\beta,\\
0, & \alpha\neq \beta,
\end{cases}
\]
which yields the displayed block decomposition and shows that the center is generated by the block projectors. Within each block, observables from \(A\cup B\) act only on the left factor and observables from \(B\cup D\) act only on the right factor. If the state is exact Markov, the standard HJPW structure theorem gives the blockwise factorized normal form \cite{HJPW}. The same formula is used in the idealized recoverability limit when exact identities are taken literally.
\end{proof}

\begin{remark}[Status of the regulator package and remaining obstruction]
Proposition~\ref{prop:regulatedrealization} and Theorem~\ref{thm:markovcollar} show that, on the ordinary or central-defect branch, edge-center completion is a theorem of overlap consistency plus the finite regulator presentation; no preferred quantum-link realization is being assumed there. Assumption~\ref{ass:technical}\textbf{T1} is therefore theorem-local shorthand for this derived regulator package. The remaining fixed-cutoff obstruction to a fully general ordinary-group theorem is a genuinely noncentral 2-group defect, which would require a higher-categorical replacement for the direct-sum center above. The later open burden is the controlled refinement/scaling regime of Assumption~\ref{ass:technical}\textbf{T8} and the modular or null-limit premises, not the matrix-algebra bookkeeping at one regulator scale.
\end{remark}

\begin{proposition}[Approximate collar recovery]\label{prop:approxrecover}
Let $A$-$B$-$D$ be a collar tripartition satisfying
\[
I(A:D|B)_\rho\le \varepsilon,
\]
and let $\mathcal R_{B\to BD}$ be a recovery map such that
\[
F\!\left(\rho_{ABD},(\mathrm{id}_A\otimes \mathcal R_{B\to BD})(\rho_{AB})\right)\ge e^{-\varepsilon/2}.
\]
Define the recovered state
\[
\rho^{\mathrm{rec}}_{ABD}:=(\mathrm{id}_A\otimes \mathcal R_{B\to BD})(\rho_{AB}).
\]
Then every bounded observable $X$ supported on $A\cup B$ obeys
\[
\left|\mathrm{Tr}\!\left[X(\rho_{ABD}-\rho^{\mathrm{rec}}_{ABD})\right]\right|
\le
\|X\|_\infty\,\|\rho_{ABD}-\rho^{\mathrm{rec}}_{ABD}\|_1
\le
2\|X\|_\infty \sqrt{1-e^{-\varepsilon}}.
\]
\end{proposition}

\begin{proof}
The first inequality is the trace-norm bound for bounded observables. The second is the Fuchs--van de Graaf estimate applied to the fidelity lower bound supplied by recoverability \cite{FawziRenner}. Thus small conditional mutual information controls deviations from the exact Markov splice identities by $O(\varepsilon^{1/2})$.
\end{proof}

\begin{remark}
Theorems~\ref{thm:markovcollar} and Corollary~\ref{cor:n1} below use exact Markovity or an explicitly stated idealized recoverability limit. Axiom~\ref{ax:entropy} itself supplies only approximate recoverability. The compact note therefore separates exact identities from their controlled $O(\varepsilon^{1/2})$ approximations rather than silently upgrading one regime into the other.
\end{remark}

\begin{proposition}[Exact Markov-collar splice theorem]\label{prop:splice}
Under Theorem~\ref{thm:markovcollar}, let $\sigma^{(\alpha)}_{b_R^\alpha D'}$ be any normalized family of states compatible with the same right-boundary sectors, and define
\[
\rho'_{AB D'}
=
\bigoplus_\alpha p_\alpha\,
\rho^{(\alpha)}_{A b_L^\alpha}\otimes \sigma^{(\alpha)}_{b_R^\alpha D'}.
\]
Define the common left algebra
\[
\mathcal A_{A b_L}:=
\bigoplus_\alpha \mathcal B(\mathcal H_{A b_L^\alpha})\otimes \mathbf 1_{b_R^\alpha},
\]
canonically represented on both direct-sum Hilbert spaces. Then every $X\in\mathcal A_{A b_L}$ has the same expectation value in $\rho_{ABD}$ and $\rho'_{AB D'}$:
\[
\mathrm{Tr}(X\rho'_{AB D'})=\mathrm{Tr}(X\rho_{ABD}).
\]
\end{proposition}

\begin{proof}
Blockwise factorization gives
\[
\mathrm{Tr}(X\rho'_{AB D'})
=
\sum_\alpha p_\alpha
\mathrm{Tr}\!\left(X_\alpha\rho^{(\alpha)}_{A b_L^\alpha}\right)
\mathrm{Tr}\!\left(\sigma^{(\alpha)}_{b_R^\alpha D'}\right).
\]
Here $X_\alpha$ denotes the $\alpha$-block of $X$ in $\mathcal A_{A b_L}$.
Each $\sigma^{(\alpha)}$ is normalized, so the second factor is $1$. The same computation for $\rho_{ABD}$ gives the same value because the original right factor is likewise normalized.
\end{proof}

\begin{proposition}[Generalized-entropy split]\label{prop:gensplit}
If the reduced cap state takes the block form
\[
\rho_C=
\bigoplus_\alpha p_\alpha
\left(
\rho^{(\alpha)}_{\mathrm{bulk},C}\otimes
\frac{\mathbf 1^{(\alpha)}_{\mathrm{edge}}}{d_\alpha}
\right),
\]
then
\[
S(\rho_C)=S_{\mathrm{bulk}}(C)+\mathrm{Tr}(\rho_C L_C),
\]
where
\[
S_{\mathrm{bulk}}(C)=H(p_\alpha)+\sum_\alpha p_\alpha S(\rho^{(\alpha)}_{\mathrm{bulk},C}),
\qquad
L_C=\sum_\alpha (\log d_\alpha)P_\alpha.
\]
\end{proposition}

\begin{proof}
The entropy of a direct sum is
\[
S\!\left(\bigoplus_\alpha p_\alpha \sigma_\alpha\right)
=
H(p_\alpha)+\sum_\alpha p_\alpha S(\sigma_\alpha).
\]
Apply this identity with
\[
\sigma_\alpha=
\rho^{(\alpha)}_{\mathrm{bulk},C}\otimes
\frac{\mathbf 1^{(\alpha)}_{\mathrm{edge}}}{d_\alpha}.
\]
Since
\[
S(\sigma_\alpha)=S(\rho^{(\alpha)}_{\mathrm{bulk},C})+\log d_\alpha,
\]
the first statement follows.
\end{proof}

\begin{definition}[Effective Newton constant in the refinement dictionary]\label{def:newton}
In the refinement-scaling regime, if
\[
\mathrm{Tr}(\rho_C L_C)\approx N_\Sigma \bar\ell(t),
\qquad
A(\partial C)\approx N_\Sigma a_{\mathrm{cell}},
\]
then matching $\mathrm{Tr}(\rho_C L_C)$ to the area term $A(\partial C)/(4G)$ defines the effective Newton constant
\[
G:=\frac{a_{\mathrm{cell}}}{4\bar\ell(t)}.
\]
This is a continuum dictionary identification, not part of the exact algebraic proposition.
\end{definition}

Proposition~\ref{prop:splice} is the algebraic version of interior invariance under compatible exterior substitutions. Proposition~\ref{prop:gensplit} isolates the exact entropy split, while Definition~\ref{def:newton} supplies the continuum identification that turns the edge center into a gravitational coupling.

\section{Modular Geometry, Relativity, and Einstein Dynamics}\label{sec:gravity}

\subsection{Geometric modular flow}

The continuum interpretation replaces the discrete patch net by cap algebras on $S^2$. Recoverability localizes the modular defect to an arbitrarily thin collar around the entangling surface.

\begin{theorem}[Conditional classification of cap modular flow in the geometric phase]\label{thm:bw}
Assume Axioms~\ref{ax:screen}--\ref{ax:entropy} and Assumption~\ref{ass:technical}, including cap isotropy / geometric modular action around the entangling circle and Euclidean regularity. Then in the refinement/scaling limit of the cap net, for each cap $C\subset S^2$, the modular flow is the unique cap-preserving conformal dilation with KMS normalization $2\pi$:
\[
K_C=2\pi B_C.
\]
\end{theorem}

\begin{proof}[Proof sketch]
Markov locality restricts the generator to the collar in the refinement/scaling regime. The assumed geometric modular phase together with cap isotropy leaves a unique noncompact one-parameter subgroup compatible with the cap. Euclidean regularity fixes the normalization of that subgroup to $2\pi$, as in the Bisognano-Wichmann modular geometry of wedge algebras \cite{BW75,BW76}.
\end{proof}

\begin{corollary}[Lorentz Kinematics]\label{cor:lorentz}
Under the hypotheses of Theorem~\ref{thm:bw},
\[
\Conf^+(S^2)\cong \mathrm{PSL}(2,\mathbb C)\cong \mathrm{SO}^+(3,1).
\]
The induced local kinematic group is the connected Lorentz group.
\end{corollary}

\subsection{The null modular bridge}

The gravity chain requires a bridge from modular data to a local stress tensor. In the compact note this bridge is pursued within the OPH framework, but the final null-stress identification is still a scaling-limit ingredient tied to Assumption~\ref{ass:technical}; it is not presented as a fixed-cutoff matrix-algebra theorem.

For a small entangling circle $\Sigma=\partial C$, let $\mathcal N$ be the null surface generated by orthogonal null geodesics. Label generators by angle $\Omega$ and affine parameter $v$.

The derivation has five steps:
\begin{enumerate}[leftmargin=2em]
\item null-strip edge-center completion;
\item modular additivity on consecutive null intervals;
\item weak continuity and finite variation of interval modular Hamiltonians;
\item emergence of a positive null translation generator from half-sided modular inclusion;
\item scaling-limit identification of the renormalized modular Hamiltonian with a null stress-tensor charge.
\end{enumerate}
Only the last step uses continuum stress-tensor language. The previous steps refer directly to the overlap algebra and its regulated collar structure.

\begin{proposition}[Null Edge-Center Completion]\label{prop:nullec}
Under the regulator premises in Assumption~\ref{ass:technical}, a regulated null buffer strip has the same central block decomposition as an ordinary spatial collar. If the reference/MaxEnt state on consecutive strips satisfies exact Markovity
\[
I_\omega(I_-:I_+|J)=0,
\]
or the explicitly assumed idealized recoverability limit that reduces to it, then the strip modular data satisfy the same state-level factorization used below.
\end{proposition}

\begin{proof}[Proof sketch]
At fixed regulator scale, a cut at constant affine parameter $v$ is a boundary for the strip region in exactly the same kinematic sense in which $\partial C$ is a boundary for a spatial collar. Type-I factorization at the regulator and boundary gauge invariance therefore imply a block decomposition
\[
\mathcal H_J \cong \bigoplus_{\alpha_1,\alpha_2}
\mathcal H_{j_L^{\alpha_1,\alpha_2}}\otimes
\mathcal H_{j_R^{\alpha_1,\alpha_2}},
\]
with the center generated by the block projectors. Within each block, observables in $I_-\cup J$ act only on the left factor and observables in $J\cup I_+$ act only on the right factor. This is the null analogue of the ordinary edge-center completion argument. The additional Markov statement is a condition on the reference state, not a consequence of the block decomposition alone.
\end{proof}

\begin{corollary}[Null Modular Additivity]\label{cor:n1}
Under Proposition~\ref{prop:nullec} and the exact Markovity condition
\[
I_\omega(I_-:I_+|J)=0,
\]
\[
K_{I_-\cup J\cup I_+}=K_{I_-\cup J}+K_{J\cup I_+}-K_J
\]
up to central boundary terms.
\end{corollary}

\begin{proof}
For a Markov state, the modular Hamiltonian satisfies the additive identity above blockwise. The boundary term is central because it depends only on the strip-edge labels.
\end{proof}

\begin{proposition}[Weak Continuity of Null Modular Hamiltonians]\label{prop:n3}
Under Axiom~\ref{ax:maxent}, Assumption~\ref{ass:technical}\textbf{T2}, and Assumption~\ref{ass:technical}\textbf{T3}, matrix elements of $K[I]$ obey the interval-length domination estimate
\[
\bigl|\langle \psi,(K[I']-K[I])\phi\rangle\bigr|
\le C_{\psi,\phi}\,\lvert I'\Delta I\rvert
\]
on bounded intervals. In particular they vary continuously with interval endpoints and have finite variation there.
\end{proposition}

\begin{proof}[Proof sketch]
Local MaxEnt gives a quasi-local Gibbs form for the regulated state. Quasi-locality together with finite-velocity propagation implies that moving an interval endpoint by $\Delta v$ changes expectation values of local observables by $O(\Delta v)$ after smearing, while Assumption~\ref{ass:technical}\textbf{T3} supplies the required endpoint-Lipschitz control. The resulting matrix elements are therefore Lipschitz on bounded intervals and in particular have finite variation.
\end{proof}

\begin{definition}[Renormalized null modular functional]
For a null interval $I$, write
\[
\widetilde K[I]:=K[I]-K_\partial(I)
\]
for the full modular Hamiltonian with the central endpoint-label term removed.
\end{definition}

\begin{theorem}[Half-line endpoint-differentiation formula]\label{lem:density}
Fix a null generator $\Omega$ and define
\[
\widetilde K_a(\Omega):=\widetilde K[(a,\infty),\Omega].
\]
Assume that for fixed vectors $\psi,\phi$ in a common dense domain the function
\[
f_{\psi,\phi}(a):=\langle\psi,\widetilde K_a(\Omega)\phi\rangle
\]
is twice weakly differentiable in $a$ with
\[
\lim_{a\to+\infty}f_{\psi,\phi}(a)=0,
\qquad
\lim_{a\to+\infty}\partial_a f_{\psi,\phi}(a)=0.
\]
Then there exists a sesquilinear-form density $p(v,\Omega)$ such that
\[
\langle\psi,\widetilde K_a(\Omega)\phi\rangle
=
2\pi\int_a^\infty (v-a)\,\langle\psi,p(v,\Omega)\phi\rangle\,dv,
\]
and distributionally
\[
\langle\psi,p(a,\Omega)\phi\rangle
=
\frac{1}{2\pi}\,\partial_a^2 f_{\psi,\phi}(a).
\]
\end{theorem}

\begin{proof}
Define the matrix elements of $p$ distributionally by
\[
\langle\psi,p(a,\Omega)\phi\rangle
:=
\frac{1}{2\pi}\,\partial_a^2 f_{\psi,\phi}(a).
\]
Integrating twice from $a$ to $+\infty$ and using the assumed vanishing of $f_{\psi,\phi}$ and its first derivative at infinity yields
\[
f_{\psi,\phi}(a)
=
2\pi\int_a^\infty (v-a)\,\langle\psi,p(v,\Omega)\phi\rangle\,dv.
\]
This is the desired half-line formula. Unlike a sigma-additive bounded-interval measure statement, it is compatible with later interval-dependent kernels on finite null intervals.
\end{proof}

\begin{corollary}[Half-Sided Modular Inclusion]\label{cor:n2}
Assume in addition the standard half-sided modular inclusion premise for the null half-line blow-up algebras,
\[
\sigma_t^\omega(\mathcal A(v>v_1))\subseteq \mathcal A(v>v_1)\qquad (t\ge 0).
\]
Then the blow-up limit of Theorem~\ref{thm:bw} on a smooth entangling cut yields a null dilation
\[
v-v_0 \mapsto e^{-2\pi t}(v-v_0),
\]
which is the geometric manifestation of the assumed half-sided modular inclusion for null half-lines.
\end{corollary}

\begin{proof}
Near a smooth cut, the cap dilation from Theorem~\ref{thm:bw} becomes a Rindler boost in the tangent geometry. Restricting that boost to a null generator gives the stated dilation. For $t\ge 0$, the half-line $v>v_0$ maps into itself, matching the stated half-sided inclusion property.
\end{proof}

\begin{lemma}[Null Translation Generator]\label{lem:translation}
Under the standard half-sided modular inclusion of Corollary~\ref{cor:n2}, there exists a one-parameter unitary group $U(a)=e^{iaP}$ with $P\ge 0$ such that
\[
\Delta^{it}U(a)\Delta^{-it}=U(e^{-2\pi t}a),
\qquad
[K,P]=i2\pi P.
\]
\end{lemma}

\begin{proof}[Proof sketch]
This is the Borchers-Wiesbrock consequence of half-sided modular inclusion. The modular group of the larger algebra acts by endomorphisms on the smaller one for one sign of modular time, and that structure produces a positive translation generator $P$ whose affine scaling law is the one written above.
\end{proof}

\begin{theorem}[Conditional null-stress identification from modular data]\label{thm:stress}
Under Axioms~\ref{ax:screen}--\ref{ax:entropy}, Assumption~\ref{ass:technical}, the hypotheses of Theorem~\ref{lem:density}, and the scaling-limit null-stress identification premise of the relativistic EFT branch, the renormalized modular Hamiltonian on a null half-line or bounded interval takes the geometric form
\[
\widetilde K[(a,\infty),\Omega]
=
2\pi\int_a^\infty (v-a)\,T_{kk}(v,\Omega)\,dv
+E_{(a,\infty),\Omega}^{(\varepsilon)},
\]
with
\[
\bigl|\langle\psi,E_{I,\Omega}^{(\varepsilon)}\phi\rangle\bigr|
\le
C_{I,\Omega,\psi,\phi}\,\varepsilon^{1/2}
\]
for all domain vectors $\psi,\phi$, where $\varepsilon$ is the recoverability/scaling-limit control parameter. If bounded null intervals inherit geometric modular action from the interval-preserving projective flow, then for each bounded interval $I$ one likewise has
\[
\widetilde K[I,\Omega]
=
2\pi\int_I g_I(v)\,T_{kk}(v,\Omega)\,dv
+E_{I,\Omega}^{(\varepsilon)},
\]
where $g_I(v)$ is affine-covariant and vanishes at finite interval endpoints. The full modular Hamiltonian is then
\[
K[I,\Omega]=\widetilde K[I,\Omega]+K_\partial(I,\Omega),
\]
with $K_\partial$ central/boundary-supported. Knowledge of $T_{kk}$ for all null directions reconstructs a symmetric tensor $T_{ab}$ up to the addition of a metric term $\phi g_{ab}$.
\end{theorem}

\begin{proof}[Proof sketch]
Theorem~\ref{lem:density} gives the half-line endpoint-differentiation formula for the renormalized modular family. Corollary~\ref{cor:n2} and Lemma~\ref{lem:translation} identify the resulting half-line density with the positive generator of null translations. In the relativistic scaling-limit branch, one then uses the standard null modular-Hamiltonian identification for local stress-tensor charges, while the recoverability/scaling-limit remainder is absorbed into the error operator $E_{I,\Omega}^{(\varepsilon)}$. If the bounded-interval modular action is geometric, the same half-line density transports to finite intervals with an interval-dependent kernel $g_I(v)$ rather than an interval-independent additive measure density. The ambiguity
\[
T_{ab}\to T_{ab}+\phi g_{ab}
\]
is invisible on null contractions because $g_{ab}k^ak^b=0$.
\end{proof}

\begin{lemma}[Null Data Determine the Stress Tensor Modulo a Metric Term]\label{lem:nullreconstruct}
Let $X_{ab}$ be a symmetric tensor. If
\[
X_{ab}k^ak^b=0
\]
for every null vector $k$ at a point, then
\[
X_{ab}=\phi g_{ab}
\]
for some scalar $\phi$.
\end{lemma}

\begin{proof}
Write
\[
X_{ab}=Y_{ab}+\phi g_{ab}
\]
with $Y_{ab}$ traceless. Since $g_{ab}k^ak^b=0$ on null vectors, the hypothesis implies $Y_{ab}k^ak^b=0$ for all null $k$. In local inertial coordinates let $k=(1,\hat n)$ with $|\hat n|=1$. Then
\[
Y_{ab}k^ak^b
=
Y_{00}+2\hat n^iY_{0i}+\hat n^i\hat n^jY_{ij}.
\]
Oddness under $\hat n\mapsto-\hat n$ forces $Y_{0i}=0$. Averaging the remaining relation over $S^2$ gives
\[
Y_{00}+\frac13\delta^{ij}Y_{ij}=0.
\]
Tracelessness of $Y_{ab}$ gives
\[
-Y_{00}+\delta^{ij}Y_{ij}=0,
\]
so both $Y_{00}$ and the spatial trace $\delta^{ij}Y_{ij}$ vanish. The remaining condition is therefore
\[
\hat n^i\hat n^j Y_{ij}^{\mathrm{TF}}=0
\qquad
\text{for all }\hat n\in S^2,
\]
where $Y_{ij}^{\mathrm{TF}}$ is the traceless spatial part. The $\ell=2$ spherical harmonics are linearly independent, so this implies $Y_{ij}^{\mathrm{TF}}=0$. Hence $Y_{ab}=0$ and therefore $X_{ab}=\phi g_{ab}$.
\end{proof}

\begin{remark}
Lemma~\ref{lem:nullreconstruct} explains why the Einstein equation derived from null modular data is determined only up to $\Lambda g_{ab}$. The missing scalar piece is not an error term. It is the unique null-invisible component of a symmetric tensor.
\end{remark}

\subsection{Generalized entropy and the Einstein equation}

From this point through Corollary~\ref{cor:einstein} we specialize to the physical $d=4$ scaling branch.

For a reference state $\omega$ and a small cap $C$,
\[
\delta S_C=\delta \langle K_C\rangle.
\]
By Theorem~\ref{thm:bw},
\[
\delta S_C = 2\pi\,\delta\langle B_C\rangle.
\]
\noindent
By Proposition~\ref{prop:gensplit},
\[
S_{\mathrm{gen}}(C)=\mathrm{Tr}(\rho L_C)+S_{\mathrm{bulk}}(C).
\]
Together with Definition~\ref{def:newton}, this identifies the area term as the coarse-grained form of the edge entropy density rather than an independent postulate.

\begin{definition}[Admissible equilibrium variation]
Throughout this subsection, $\delta$ denotes a first-order state variation about a reference state $\omega$, within the cap-label-preserving MaxEnt class, at fixed cap $C$, fixed cap size/volume, and fixed conserved charges. Shape or null-cut deformations are denoted separately by $\delta_{\mathrm{shape}}$ or $\partial_\lambda$.
\end{definition}

\begin{assumption}[Fixed-cap generalized-entropy stationarity]\label{ass:equilibrium}
For admissible variations preserving cap labels, cap size/volume, and relevant conserved charges, the realized reference state is stationary for generalized entropy:
\[
\delta S_{\mathrm{gen}}(C)=0.
\]
\end{assumption}

This is the Jacobson-type fixed-cap stationarity premise used in the compact note. MaxEnt alone does not derive it here.

\begin{lemma}[d=4 small-ball modular bridge]\label{lem:smallball}
Assume the hypotheses of Theorem~\ref{thm:stress} and a locally Lorentzian scaling regime in which $\delta\langle T_{00}\rangle$ is approximately constant across a small radius-$\ell$ ball/diamond centered at the cap point. Then in $d=4$,
\[
\delta S_C
=
2\pi\,\delta\langle B_C\rangle
=
2\pi\int_{B_\ell}\frac{\ell^2-r^2}{2\ell}\,\delta\langle T_{00}\rangle\,d^3x
=
\frac{8\pi^2\ell^4}{15}\,\delta\langle T_{00}\rangle
+O(\ell^5\partial T).
\]
\end{lemma}

\begin{proof}[Proof sketch]
Theorem~\ref{thm:stress} supplies the null modular/stress-tensor identification in the scaling branch. In the locally Lorentzian $d=4$ small-ball limit, the cap-preserving conformal generator is the standard ball modular kernel $(\ell^2-r^2)/(2\ell)$. Integrating that kernel over the three-ball gives
\[
\int_{B_\ell}\frac{\ell^2-r^2}{2\ell}\,d^3x=\frac{4\pi\ell^4}{15},
\]
and multiplying by the overall $2\pi$ factor from the first law yields the stated coefficient.
\end{proof}

\begin{theorem}[Conditional Jacobson-type Einstein relation from entanglement equilibrium]\label{thm:einstein}
Assume Axioms~\ref{ax:screen}--\ref{ax:entropy}, Assumption~\ref{ass:technical}, Assumption~\ref{ass:equilibrium}, and the hypotheses of Theorem~\ref{thm:stress} and Lemma~\ref{lem:smallball}. Then, for every sufficiently small cap $C$ and every admissible first variation $\delta$ about a reference state $\omega$ in the locally Lorentzian $d=4$ scaling regime,
\[
\delta\!\left(G_{00}+\Lambda g_{00}\right)=8\pi G\,\delta\langle T_{00}\rangle
\]
at the cap center in the diamond rest frame.
\end{theorem}

\begin{proof}[Proof sketch]
Assumption~\ref{ass:equilibrium} gives the fixed-cap stationarity condition
\[
\delta S_{\mathrm{gen}}(C)=0.
\]
Lemma~\ref{lem:smallball} identifies the bulk modular contribution in the $d=4$ small-ball regime as
\[
\delta S_{\mathrm{bulk}}(C)=\frac{8\pi^2\ell^4}{15}\,\delta\langle T_{00}\rangle.
\]
For a fixed-volume small ball in $d=4$, the standard area variation is
\[
\delta A\big|_{V,\Lambda_{\mathrm{ref}}}
=
-\frac{4\pi\ell^4}{15}\,
\delta\!\left(G_{00}+\Lambda_{\mathrm{ref}}g_{00}\right).
\]
Since $S_{\mathrm{gen}}=S_{\mathrm{bulk}}+A/(4G)$ in the coarse-grained scaling dictionary, stationarity gives
\[
0=\delta S_{\mathrm{bulk}}+\frac{\delta A}{4G}
=
\frac{8\pi^2\ell^4}{15}\,\delta\langle T_{00}\rangle
-\frac{\pi\ell^4}{15G}\,
\delta\!\left(G_{00}+\Lambda_{\mathrm{ref}}g_{00}\right),
\]
which rearranges to the stated rest-frame relation.
\end{proof}

\begin{corollary}[Tensor Upgrade in the Scaling Regime]\label{cor:einstein}
If the first-variation relation of Theorem~\ref{thm:einstein} holds for all local directions and reference states in the scaling regime, then
\[
G_{ab}+\Lambda g_{ab}=8\pi G\,\langle T_{ab}\rangle
\]
modulo the metric-term ambiguity already isolated by Lemma~\ref{lem:nullreconstruct}.
\end{corollary}

\begin{proof}
Overlap consistency across observers through the same point supplies all timelike directions, and Theorem~\ref{thm:stress} together with Lemma~\ref{lem:nullreconstruct} upgrades the rest-frame relation to the full tensor equation in the Jacobson entanglement-equilibrium line of argument \cite{Jacobson95,Jacobson16}.
\end{proof}

\subsection{The cosmological constant}

\begin{lemma}[Vacuum-Energy Blindness]\label{lem:vac}
For any null vector $k$ and any vacuum-energy contribution $T^{\mathrm{vac}}_{ab}=-\rho_{\mathrm{vac}}g_{ab}$,
\[
T^{\mathrm{vac}}_{kk}=0.
\]
\end{lemma}

\begin{proposition}[Structural Separation of $\Lambda$]\label{prop:lambda}
Local null-modular data fix $T_{ab}$ only up to $\phi g_{ab}$. Therefore $\Lambda$ is not determined by local overlap consistency and requires a global input.
\end{proposition}

\begin{corollary}[Cosmological Constant from Capacity]\label{cor:lambda}
If the global screen capacity is identified with the de Sitter static-patch entropy,
\[
N_{\mathrm{scr}}=S_{\mathrm{dS}}=\frac{A_{\mathrm{dS}}}{4G}=\frac{3\pi}{G\Lambda},
\]
then
\[
\Lambda=\frac{3\pi}{G N_{\mathrm{scr}}}.
\]
\end{corollary}

This is a structural aspect of the framework. The vacuum-energy problem is reorganized by separating local null data from the global capacity term.

\section{Gauge Reconstruction and Standard Model Structure}

\subsection{Compact gauge reconstruction}

At any fixed UV cutoff, edge-center completion equips collars with finitely many sector labels and fusion rules. The category used for reconstruction is the refinement-stable directed colimit
\[
\mathsf{Sect}_\infty := \varinjlim_r \mathsf{Sect}_r
\]
over regulator refinements $r$, retaining only transportable sectors and intertwiners that persist under refinement. This is the object on which Tannaka/Doplicher-Roberts reconstruction is applied.

\begin{theorem}[Conditional compact gauge reconstruction in the bosonic branch]\label{thm:compactgauge}
Assume Axioms~\ref{ax:screen}--\ref{ax:entropy}, the relevant parts of Assumption~\ref{ass:technical}, and that the refinement-stable edge-sector category $\mathsf{Sect}_\infty$ in the EFT regime is a rigid symmetric $C^*$-tensor category with a faithful bosonic fiber functor to finite-dimensional Hilbert spaces. Then
\[
\mathsf{Sect}_\infty\simeq \mathrm{Rep}(G)
\]
for a compact group $G$ uniquely determined up to isomorphism.
\end{theorem}

\begin{proof}[Proof sketch]
Each fixed-cutoff category has only finitely many simple blocks, but the refinement-stable colimit $\mathsf{Sect}_\infty$ can have infinitely many simple objects while still carrying objectwise finite-dimensional fibers. In the bosonic EFT branch, symmetry and rigidity place the category in the Doplicher-Roberts/Tannakian reconstruction class, and the faithful bosonic fiber functor identifies it with $\mathrm{Rep}(G)$ for a compact group $G$ \cite{DR89,DR90}.
\end{proof}

If a fermionic sign object is present, the correct reconstruction statement is super-Tannakian rather than purely Tannakian. The compact note does not prove that fermionic/chiral extension; it works only in the bosonic internal-gauge branch.

\subsection{Minimal admissibility and the Standard Model quotient}

Throughout this subsection, $\chi_{\mathrm{cpl}}$ means the dimension of the minimal coupled carrier supporting the required weak-type and color-type nonabelian charges on a common block. It is not the abstract minimal faithful representation dimension of the final gauge group.

\begin{definition}[Weak-type and color-type nonabelian roles]
A weak-type nonabelian charge is a pseudoreal nonabelian doublet role. A color-type nonabelian charge is an intrinsically complex nonabelian role, meaning that its nonabelian factor itself acts by an irreducible representation not equivalent to its conjugate. Abelian twisting of a pseudoreal nonabelian irreducible representation does not count as a color-type role.
\end{definition}

The derivation of the Standard Model quotient separates into five steps:
\begin{enumerate}[leftmargin=2em]
\item existence of a compact reconstructed group;
\item identification of the minimal nonabelian sector content;
\item identification of the minimal coupled carrier;
\item uniqueness of the connected abelian factor once admitted on that carrier;
\item determination of product gauge structure up to finite quotient, then of the global finite quotient from the realized matter spectrum.
\end{enumerate}

\begin{lemma}[Minimal nonabelian sector content]\label{lem:mincontent}
Any admissible low-energy sector must contain both a weak-type pseudoreal nonabelian charge role and an intrinsically complex color-type nonabelian charge role.
\end{lemma}

\begin{proof}
Light chiral matter requires left-handed multiplets and right-handed singlets to carry inequivalent gauge data. A pseudoreal doublet structure is the minimal way to realize the weak sector without immediately producing vectorlike masses. A genuinely complex nonabelian representation is required to distinguish quarks from antiquarks. A single nonabelian simple factor cannot simultaneously supply both a minimal weak-type role and an intrinsically complex color-type role; abelian twisting does not count toward the color-type requirement.
\end{proof}

\begin{lemma}[Connected 2D pseudoreal image classification]\label{lem:su2class}
Let $H$ be a connected compact group with a faithful irreducible $2$-dimensional pseudoreal unitary representation. Then the connected derived subgroup of its image is conjugate to $\SU(2)$. Equivalently, the nonabelian factor acting in that role is $\SU(2)$ up to finite central quotient.
\end{lemma}

\begin{proof}[Proof sketch]
Irreducibility places the image inside $\U(2)$ with scalar commutant. Its connected derived subgroup is therefore a connected compact subgroup of $\SU(2)$. The connected compact subgroups of $\SU(2)$ are either tori or $\SU(2)$ itself. The representation is nonabelian and pseudoreal, so the torus case is excluded. Hence the derived subgroup is conjugate to $\SU(2)$, with only a finite central kernel left invisible in the representation.
\end{proof}

\begin{lemma}[Connected 3D intrinsically complex image classification]\label{lem:su3class}
Let $H$ be a connected compact group with a faithful irreducible intrinsically complex $3$-dimensional unitary representation. Then the connected derived subgroup of its image is conjugate to $\SU(3)$. Equivalently, the nonabelian factor acting in that role is $\SU(3)$ up to finite central quotient.
\end{lemma}

\begin{proof}[Proof sketch]
Because the representation is irreducible on $\mathbb C^3$, the semisimple part of the Lie algebra of the image acts irreducibly on $\mathbb C^3$. A product of two nontrivial simple factors would force a tensor-product decomposition of dimension at least $4$, so only one simple factor can occur. Among compact simple Lie algebras, the only ones with nontrivial irreducible representations of dimension at most $3$ are $\mathfrak{su}(2)$ and $\mathfrak{su}(3)$. The $3$-dimensional $\mathfrak{su}(2)$ representation is real, not intrinsically complex, whereas the fundamental $\mathfrak{su}(3)$ representation is intrinsically complex. Therefore the connected derived subgroup is conjugate to $\SU(3)$, again up to finite central kernel.
\end{proof}

\begin{lemma}[Minimal coupled carrier under MAR$^+$]\label{lem:carrier}
Within the connected positive-dimensional Lie admissible class used by MAR, and under the weak-type / intrinsically complex color-type criteria above, the minimal coupled carrier supporting both roles on a common block has the form
\[
V=\mathbb C^3\otimes \mathbb C^2,
\qquad
\chi_{\mathrm{cpl}}=6.
\]
\end{lemma}

\begin{proof}
By Lemma~\ref{lem:su2class}, the minimal weak-type nonabelian role is the pseudoreal doublet of $\SU(2)$, hence has dimension $2$. By Lemma~\ref{lem:su3class}, the minimal intrinsically complex color-type nonabelian role is the triplet of $\SU(3)$, hence has dimension $3$.

By definition of $\chi_{\mathrm{cpl}}$, the weak-type and color-type roles must act nontrivially on one common irreducible nonabelian block rather than on disconnected direct-sum summands. Commuting irreducible actions on such a common block therefore require dimension at least the product of the minimal weak and color dimensions:
\[
\chi_{\mathrm{cpl}}\ge 2\cdot 3=6.
\]
Equality is achieved by the tensor-product carrier $\mathbb C^3\otimes\mathbb C^2$. The block-diagonal representation $\mathbb C^3\oplus \mathbb C^2$ of $S(U(3)\times U(2))$ is faithful of dimension $5$, but it is not coupled and hence is not the MAR minimizer. Connected non-semisimple counterexamples of $U(2)$ type are excluded by the explicit color-type criterion rather than being silently ignored.
\end{proof}

\begin{remark}\label{rem:connectedlie}
The dimension comparisons used in Lemma~\ref{lem:carrier} apply to the positive-dimensional connected Lie image carrying the admissible nonabelian charges, equivalently to the identity component of the reconstructed group on that sector. Finite, disconnected, or connected non-semisimple counterexamples exist, but they are outside the admissible class fixed by the explicit weak-type / color-type criteria.
\end{remark}

MAR supplies existence of a connected abelian factor; the next lemma proves only uniqueness and identification on the minimal coupled carrier.

\begin{lemma}[Uniqueness of the connected abelian factor on the minimal coupled carrier]\label{lem:commutant}
Inside $\U(6)$, the commutant of $\SU(3)\times\SU(2)$ acting on $\mathbb C^3\otimes \mathbb C^2$ is exactly $\U(1)$.
\end{lemma}

\begin{proof}
The $\SU(3)$ action is irreducible on the first tensor factor and trivial on the second. The $\SU(2)$ action is irreducible on the second tensor factor and trivial on the first. By Schur's lemma, any operator commuting with both actions is scalar on each irreducible factor. Hence the full commutant is a single $\U(1)$.
\end{proof}

\begin{theorem}[Product Gauge Structure up to Finite Quotient]\label{thm:sm}
Assume Axiom~\ref{ax:mar}, vanishing obstruction $[z]=0$ on the realized low-energy branch, the hypotheses of Theorem~\ref{thm:compactgauge}, the weak-type / color-type MAR premises above, and that the realized connected gauge image acts faithfully on the minimal coupled carrier. Then the realized connected gauge structure has the form
\[
G_{\mathrm{phys}}=
\frac{\SU(3)\times\SU(2)\times\U(1)}{\Gamma}
\]
for some finite central subgroup $\Gamma$.
\end{theorem}

\begin{proof}
Theorem~\ref{thm:compactgauge} yields some compact group $G$. Lemma~\ref{lem:mincontent} identifies the minimal admissible sector content, Lemmas~\ref{lem:su2class} and \ref{lem:su3class} identify the connected nonabelian factors realizing the minimal weak-type and color-type roles, and Lemma~\ref{lem:carrier} fixes the minimal coupled carrier. The tensor-product structure of that carrier implies commuting weak and color actions, so the connected semisimple part of the realized image is locally $\SU(3)\times\SU(2)$. Axiom~\ref{ax:mar} supplies the existence of one connected abelian charge factor acting nontrivially on the coupled carrier, and Lemma~\ref{lem:commutant} shows that any such connected abelian factor is necessarily a single $\U(1)$. A compact connected Lie group with Lie algebra $\mathfrak{su}(3)\oplus\mathfrak{su}(2)\oplus\mathfrak{u}(1)$ is a quotient of $\SU(3)\times\SU(2)\times\U(1)$ by a finite central subgroup. Faithfulness rules out an invisible extra torus, so the only remaining ambiguity is the finite central quotient $\Gamma$.
\end{proof}

\begin{theorem}[Hypercharge lattice on the realized matter package]\label{thm:hypercharge}
Assume gauge group $\SU(N_c)\times\SU(2)\times\U(1)_Y$, one generation of chiral matter $(Q,u^c,d^c,L,e^c)$, one Higgs doublet $H$, and Yukawa terms
\[
QHu^c,\qquad QH^\dagger d^c,\qquad LH^\dagger e^c.
\]
Then anomaly cancellation and Yukawa invariance imply
\[
Y_L=-N_cY_Q,\quad
Y_H=N_cY_Q,\quad
Y_u=-(N_c+1)Y_Q,\quad
Y_d=(N_c-1)Y_Q,\quad
Y_e=2N_cY_Q.
\]
Fixing the normalization by $Q=T_3+Y$ and $Q(\nu_L)=0$ gives
\[
Y_Q=\frac{1}{2N_c}.
\]
For $N_c=3$ one recovers the exact Standard Model lattice
\[
Y_Q=\frac16,\quad
Y_L=-\frac12,\quad
Y_u=-\frac23,\quad
Y_d=\frac13,\quad
Y_e=1,\quad
Y_H=\frac12.
\]
\end{theorem}

\begin{proof}
Yukawa invariance gives
\[
Y_u=-(Y_Q+Y_H),\qquad
Y_d=-Y_Q+Y_H,\qquad
Y_e=-Y_L+Y_H.
\]
The mixed anomalies yield
\[
N_cY_Q+Y_L=0,
\]
and
\[
2N_cY_Q+N_cY_u+N_cY_d+2Y_L+Y_e=0.
\]
Solving the first anomaly gives
\[
Y_L=-N_cY_Q.
\]
Substituting the Yukawa relations and $Y_L=-N_cY_Q$ into the mixed gravitational anomaly then yields
\[
Y_H=N_cY_Q,\qquad
Y_u=-(N_c+1)Y_Q,\qquad
Y_d=(N_c-1)Y_Q,\qquad
Y_e=2N_cY_Q.
\]
With these relations in place, the $\SU(N_c)^2\U(1)$ and $\U(1)^3$ anomalies cancel identically. The normalization of $Y_Q$ is fixed by the electric charge operator $Q=T_3+Y$ together with $Q(\nu_L)=0$, which implies
\[
Y_L=-\frac12=-N_cY_Q,
\qquad\text{hence}\qquad
Y_Q=\frac{1}{2N_c}.
\]
\end{proof}

\begin{corollary}[Three Generations on the realized MAR-admissible branch]\label{cor:ng}
Under the hypotheses of Theorem~\ref{thm:sm},
\[
N_g=3.
\]
\end{corollary}

\begin{proof}
Intrinsic CKM CP violation requires
\[
\frac{(N_g-1)(N_g-2)}{2}>0,
\]
hence $N_g\ge 3$. For one Higgs doublet, the one-loop weak-sector coefficient is
\[
b_2=\frac{11}{3}C_A-\frac{2}{3}\sum_f T(R_f)-\frac{1}{3}\sum_s T(R_s)
=
\frac{22}{3}-\frac{N_g(N_c+1)}{3}-\frac16
=
\frac{43}{6}-\frac{N_g(N_c+1)}{3}.
\]
Asymptotic freedom therefore implies
\[
N_g(N_c+1)<\frac{43}{2}.
\]
Since the color-type requirement forces $N_c\ge 3$, this gives $N_g\le 5$. The minimal admissibility selector then chooses $N_g=3$.
\end{proof}

\begin{corollary}[Three Colors on the realized MAR-admissible branch]\label{cor:nc}
Under the hypotheses of Theorem~\ref{thm:sm},
\[
N_c=3.
\]
\end{corollary}

\begin{proof}
Witten's global $\SU(2)$ anomaly requires the total number of left-handed weak doublets to be even:
\[
N_g(N_c+1)\equiv 0 \pmod 2.
\]
With $N_g=3$, this implies $N_c$ is odd. The case $N_c=1$ is inadmissible because it removes the intrinsically complex color sector demanded by Lemma~\ref{lem:mincontent}. Minimal admissibility therefore selects $N_c=3$ \cite{Witten82}.
\end{proof}

\begin{proposition}[Global quotient from the realized matter spectrum]\label{prop:z6}
If the realized matter spectrum carries the hypercharges of Theorem~\ref{thm:hypercharge} with $N_c=3$, then the subgroup of $\SU(3)\times\SU(2)\times\U(1)$ acting trivially on all realized states is exactly $\mathbb Z_6$.
\end{proposition}

\begin{proof}
Write $q:=6Y\in\mathbb Z$ and $\eta:=e^{i\pi/3}$, so the $\U(1)$ factor acts by $\eta^q$ on charge-$q$ multiplets. The element
\[
g_0:=(\omega_3,-1,\eta)
\]
acts trivially on the realized matter and Higgs multiplets:
\[
Q:\ \omega_3(-1)\eta=1,\qquad
u^c:\ \omega_3^{-1}\eta^{-4}=1,\qquad
d^c:\ \omega_3^{-1}\eta^2=1,
\]
\[
L:\ (-1)\eta^{-3}=1,\qquad
e^c:\ \eta^6=1,\qquad
H:\ (-1)\eta^3=1.
\]
Its powers therefore generate a subgroup of order six acting trivially on all realized states. Conversely, any central element acting trivially on $e^c$ must have $\U(1)$ part $\eta^n$, and triviality on $Q$ then fixes the $\SU(3)$ and $\SU(2)$ center factors uniquely. Hence every trivial central element is a power of $g_0$, so the kernel is exactly $\mathbb Z_6$.
\end{proof}

\begin{corollary}[Standard Model Gauge Group]
Under Theorem~\ref{thm:sm}, Theorem~\ref{thm:hypercharge}, Corollaries~\ref{cor:ng} and \ref{cor:nc}, and Proposition~\ref{prop:z6},
\[
G_{\mathrm{phys}}=
\frac{\SU(3)\times\SU(2)\times\U(1)}{\mathbb Z_6}.
\]
\end{corollary}

\begin{proof}
Theorem~\ref{thm:sm} yields the product structure up to finite quotient, and Proposition~\ref{prop:z6} fixes that quotient to $\mathbb Z_6$ once the realized hypercharge lattice and $N_c=3$ are in place.
\end{proof}

\begin{corollary}[Product-Group Consequences]
The adjoint representation of the full connected gauge group contains only
\[
(8,1,0)\oplus(1,3,0)\oplus(1,1,0).
\]
Equivalently, the adjoint representation of the derived gauge group contains only
\[
(8,1,0)\oplus(1,3,0).
\]
There are no mixed $(3,2,\pm 5/6)$ gauge bosons. Hence gauge-mediated proton decay is absent. The unbroken gauge symmetries forbid elementary gauge-boson mass terms: the photon remains massless, while gluons are massless gauge fields but confined.
\end{corollary}

\begin{remark}
The same product-group structure removes the usual grand-unification monopole sector. In the present framework, unification of couplings does not proceed by embedding the gauge group into a simple group that is later broken. It proceeds through a shared geometric origin of the edge heat-kernel parameters. Absence of gauge-mediated proton decay and unification therefore coexist without tension.
\end{remark}

\section{Two-Input Quantitative Sector}

\subsection{Continuous-input count in the current quantitative implementation}

The quantitative program uses the same axiom set together with the two external continuous inputs $P$ and $N_{\mathrm{scr}}$ in its current supplement-backed numerical implementation. In the current implementation, $P$ is calibrated from the gauge sector and $N_{\mathrm{scr}}$ is inferred from the observed cosmological constant. The discrete structure
\[
(\SU(3)\times\SU(2)\times\U(1))/\mathbb Z_6,\quad
N_c=3,\quad
N_g=3,\quad
\varepsilon=\frac16,\quad
\delta=\frac29
\]
is then fixed by structural arguments, MAR selection, or phenomenological continuation steps rather than by direct fits. On the structural, calibration, and Higgs/top branch alone, this replaces roughly nineteen particle-physics inputs, or twenty-six when neutrino parameters are included, by two external continuous inputs plus structural relations. Later flavor continuations use additional phenomenological branch choices and are not part of that count.

\subsection{Epistemic classification of outputs}

The quantitative sector is not uniform. The paper uses four categories.

\begin{enumerate}[leftmargin=2em]
\item \textbf{Calibration-sector consistency checks.} These are algebraically entangled with the determination of $P$ from the gauge sector. They include $v$, the gauge couplings at $m_Z$, and tree-level electroweak masses.
\item \textbf{Independent quantitative outputs.} These use additional structural conditions and add no further continuous fit once the gauge trajectory and scale-setting branch are fixed. In the compact note this category is the Higgs/top critical-surface branch.
\item \textbf{Parameter-free structural outputs.} These include the gauge quotient, the exact hypercharge lattice on the realized matter package, and $N_c=3$, $N_g=3$ on the realized MAR-admissible branch.
\item \textbf{Phenomenological continuations.} These include the charged-lepton/Koide construction, quark textures, baryogenesis, dark-sector phenomenology, and any use of the $\mathbb Z_6$ suppression factor as a flavor or entropy ansatz.
\end{enumerate}

The quark-hierarchy and neutrino branches are not at the same maturity level in the compact note: the former is a texture-pattern continuation, while the latter is currently only an order-of-magnitude capacity-branch estimate.

This classification matters because the framework is not presented as a fit to all observables. Its stronger claim is that a large set of discrete and dimensionless structures are fixed before the overall energy scale is chosen.

\subsection{Fundamental scales}

The unification scale is
\[
M_U=\frac{E_P}{e^{2\pi}}\,P^{1/6}\approx 2.474\times 10^{16}\ \mathrm{GeV},
\]
and the cell energy is
\[
E_{\mathrm{cell}}=\frac{E_P}{\sqrt P}\approx 9.56\times 10^{18}\ \mathrm{GeV}.
\]
The factor $e^{2\pi}$ comes from modular geometry. The factor $P^{1/6}$ comes from the cell-area scaling relation.

\subsection{The pixel constraint and gauge closure}

Edge-center completion and MaxEnt produce open-edge Gibbs / Laplacian weights
\[
p_R(t)\propto d_R e^{-tC_2(R)},
\qquad
\bar\ell_G(t)=\sum_R p_R(t)\log d_R.
\]
For the two nonabelian factors, these are explicitly
\[
\bar\ell_{\SU(2)}(t)=\sum_{j=0,\frac12,1,\dots} p_j(t)\log(2j+1),
\qquad
p_j(t)\propto (2j+1)e^{-t\,j(j+1)},
\]
and
\[
\bar\ell_{\SU(3)}(t)=\sum_{p,q\ge 0} p_{(p,q)}(t)\log d_{(p,q)},
\qquad
p_{(p,q)}(t)\propto d_{(p,q)}e^{-tC_2(p,q)},
\]
with
\[
d_{(p,q)}=\frac12(p+1)(q+1)(p+q+2),
\]
and
\[
C_2(p,q)=\frac13(p^2+q^2+pq+3p+3q).
\]
Upon gluing or tracing the two boundary indices one obtains the closed heat-kernel partition function
\[
Z_{\mathrm{edge}}(t)=\sum_R d_R^2 e^{-tC_2(R)}.
\]
For the nonabelian factors one obtains the pixel constraint
\[
\frac{P}{4}=\bar\ell_{\SU(2)}(t_2)+\bar\ell_{\SU(3)}(t_3),
\qquad
t_i=4\pi^2\alpha_i.
\]

The framework then solves a supplement-backed self-consistent fixed-point problem:
\begin{enumerate}[leftmargin=2em]
\item choose a trial unified coupling $\alpha_U$;
\item run the three gauge couplings to a trial scale using the derived beta-function shifts;
\item compute the electroweak scale
\[
v=E_{\mathrm{cell}}\exp\!\left(-\frac{2\pi}{\beta_{\mathrm{EW}}\alpha_U}\right),
\qquad
\beta_{\mathrm{EW}}=N_c+1=4;
\]
\item impose $m_Z(\mu)=\mu$ at the running scale;
\item solve the pixel constraint by bisection in $\alpha_U$.
\end{enumerate}

This yields
\[
\alpha_U\approx 0.04112,
\qquad
\alpha_U^{-1}\approx 24.32.
\]

The resulting beta-function shifts are numerically close to the MSSM one-loop shifts. In the framework this is not introduced by hand. It comes from the same heat-kernel edge weighting that enters the pixel constraint. The compact note quotes this closure only at the level needed to state the numerical branch; the full beta-function shifts, threshold choices, and matching conventions remain supplement-backed.

\subsection{Calibration-sector electroweak outputs}

We use the convention
\[
\alpha_1=\frac53\alpha_Y,
\qquad
g_Y^2=4\pi\alpha_Y=4\pi\frac35\alpha_1.
\]
The gauge couplings at $\mu=m_Z$ are then
\[
\alpha_1(m_Z)\approx 0.01696,\qquad
\alpha_2(m_Z)\approx 0.03384,\qquad
\alpha_3(m_Z)\approx 0.1183,
\]
which, under the printed conventions of this compact note, give
\[
\alpha_{\mathrm{em}}^{-1}(m_Z)\approx 127.82,\qquad
\sin^2\theta_W(m_Z)\approx 0.23119.
\]
With
\[
v\approx 246.77\ \mathrm{GeV},
\]
the tree-level electroweak masses are
\[
m_W=\frac12 v g_2 \approx 80.46\ \mathrm{GeV},
\]
and
\[
m_{Z,\mathrm{run}}=\frac12v\sqrt{g_2^2+g_Y^2}\approx 91.76\ \mathrm{GeV}.
\]

These are calibration-sector outputs. Their role in the paper is scale setting, not independent confirmation. A fully reproducible numerical audit requires the beta-function shifts and matching conventions supplied in the repository supplements. The compact note therefore reports only the running-$Z$ relation internally and does not treat an illustrative pole-mass estimate as a benchmark output.

\begingroup
\footnotesize
\begin{longtable}{@{}p{0.28\textwidth}p{0.18\textwidth}p{0.18\textwidth}p{0.23\textwidth}@{}}
\toprule
Quantity & OPH output & Experimental value & Status \\
\midrule
\endhead
\bottomrule
\endfoot
$\alpha_{\mathrm{em}}^{-1}(m_Z)$ & $127.82$ & $127.952$ & calibration sector \\
$\sin^2\theta_W(m_Z)$ & $0.23119$ & $0.23122$ & calibration sector \\
$\alpha_s(m_Z)$ & $0.1183$ & $0.1179$ & calibration sector \\
$m_W$ & $80.46~\mathrm{GeV}$ & $80.38~\mathrm{GeV}$ & calibration sector \\
\end{longtable}
\endgroup

\subsection{Higgs and top from the critical surface}

The independent quantitative branch begins with the refinement-stability condition
\[
\lambda(M_U)=0,
\qquad
\beta_\lambda(M_U)=0.
\]
At one loop, these conditions determine the top Yukawa boundary value
\[
y_t(M_U)=
\left[\frac{1}{16}\left(2g_2^4+(g_2^2+g_Y^2)^2\right)\right]^{1/4}.
\]
The logic is simple: once $\lambda(M_U)=0$, the condition $\beta_\lambda(M_U)=0$ removes the final freedom in the top Yukawa boundary value. The low-energy Higgs and top masses then follow by RG evolution rather than by direct fitting.

Running the coupled system $(y_t,\lambda,g_Y,g_2,g_3)$ down to $\mu\sim m_t$ yields
\[
m_t^{\overline{\mathrm{MS}}}(m_t)\approx 160.6\ \mathrm{GeV},
\]
and after the standard three-loop QCD pole-mass conversion used in the present implementation,
\[
m_t^{\mathrm{pole}}\approx 171.1\ \mathrm{GeV}.
\]
The corresponding Higgs mass is
\[
m_H=\sqrt{2\lambda(m_t)}\,v\approx 126.5\ \mathrm{GeV}.
\]

The critical-surface boundary conditions are not an additional continuous fit once the gauge trajectory and scale-setting branch are fixed. Converting the resulting dimensionless ratios $m_H/v$ and $m_t/v$ to GeV still uses the previously calibrated $v$ and gauge-coupling trajectory, and the critical-surface branch can fail even when the electroweak calibration succeeds.

\begingroup
\footnotesize
\begin{tabular}{@{}p{0.28\textwidth}p{0.18\textwidth}p{0.18\textwidth}p{0.23\textwidth}@{}}
\toprule
Quantity & OPH output & Experimental value & Status \\
\midrule
$m_H$ & $126.5~\mathrm{GeV}$ & $125.2~\mathrm{GeV}$ & independent quantitative \\
$m_t^{\mathrm{pole}}$ & $171.1~\mathrm{GeV}$ & $172.6~\mathrm{GeV}$ & independent quantitative \\
\bottomrule
\end{tabular}
\endgroup

\section{Phenomenological Flavor Continuations}

Theorem~\ref{thm:master} packages the compact note's recovered core: the D1--D5 relativity chain together with the D7--D9 Standard Model structural chain. D6 is a separate input-dependent capacity corollary, and D10--D11 are supplement-backed numerical branches of the current implementation. Nothing in the present section is part of that recovered core. The role of the section is only to record which flavor ideas have been explored and exactly why they remain outside the proved package.

\subsection{The $\mathbb Z_6$ suppression parameter}\label{subsec:z6-suppression}

If one adds the extra assumption of a \emph{uniform} sixfold center-label ensemble on the realized $\mathbb Z_6$ quotient sectors, then the associated center-label entropy is
\[
\Delta S_{\mathrm{center}}=\ln 6.
\]
and one may define the phenomenological hierarchy parameter
\[
\varepsilon=e^{-\Delta S_{\mathrm{center}}}=\frac16.
\]
The crucial scope point is that the quotient
\[
(\SU(3)\times\SU(2)\times\U(1))/\mathbb Z_6
\]
does \emph{not} by itself force the uniform ensemble. The step from the proved quotient to the numerical suppression factor $\varepsilon=1/6$ is therefore an additional flavor ansatz, not a theorem of the recovered core.

\subsection{Koide, charged-lepton, and texture branches}

Every continuation in which one introduces a circulant generation-space ansatz, a fixed Koide phase such as
\[
Q=\frac23,\qquad \delta=\frac29,
\]
or additional texture exponents depends on extra phenomenological structure beyond the recovered core.
These constructions may be interesting as model-building continuations because they reuse discrete data already fixed on the realized Standard Model branch, but they are not derived by overlap consistency, the scaling-limit gravity premises, compact gauge reconstruction, anomaly cancellation, or MAR alone.

Accordingly, the compact note does \emph{not} claim the charged-lepton masses, the Koide relation, or texture exponents as recovered outputs. At present they should be read only as conjectural downstream organizations of flavor data. A failure of any such ansatz would retract that continuation, not Theorem~\ref{thm:master}.

\subsection{Neutrino scale}

The neutrino discussion is weaker still in scope. The present note contains only an input-dependent, order-of-magnitude capacity estimate tied to $N_{\mathrm{scr}}$; it does not provide a benchmark neutrino mass formula, a mixing derivation, or a completed flavor mechanism. The estimate itself sits at the separate D6 boundary, while any fuller neutrino branch remains outside both the recovered core and the supplement-backed D10--D11 quantitative implementation.

Phase I therefore ends \emph{before} flavor is derived. The compact paper recovers spacetime dynamics and Standard Model gauge/matter structure first; flavor remains an explicitly deferred problem rather than a claim hidden inside the packaging.

\section{Program-Level Phenomenology and Further Continuations}

Theorem~\ref{thm:master} marks the end of the compact note's recovered core. The separate input-dependent capacity corollary D6 and the D10--D11 numerical branches are kept visible, but the items below go beyond that scope. They are recorded only to show where the broader OPH program might go next; they are \emph{not} theorem-level outputs of the compact paper.

\subsection{What is already fixed before any continuation is added}

Two points should be separated from the downstream branches. First, the product-group consequence
\[
\tau_p^{(\mathrm{gauge})}=\infty
\]
for gauge-mediated proton decay is already a corollary of the realized Standard Model structure and does not depend on any dark-sector, baryogenesis, string, or spectroscopy continuation. Second, the relation
\[
\Lambda=\frac{3\pi}{G N_{\mathrm{scr}}}
\]
is already isolated earlier as the separate input-dependent capacity corollary D6 once the screen-capacity identification is adopted. Neither point should be confused with the conjectural branches below.

\subsection{Dark-sector branch}

The modular-anomaly proposal for an effective dark stress tensor requires more than the recovered gravity branch. The compact note does derive the modular additivity defect
\[
\Delta K_\delta = K_{ABD}-K_{AB}-K_{BD}+K_B
\]
as a recoverability diagnostic, but it does not prove that a nontrivial cap-local anomaly
\[
K_C^{(\mathrm{anom})}:=K_C-2\pi B_C
\]
survives with the sign, magnitude, and infrared closure required to explain galaxy-scale phenomenology. Any MOND/RAR-type acceleration law extracted from this idea therefore remains a conjectural response model, not a compact-note prediction.

\subsection{Baryogenesis, spectroscopy, and string/worldsheet branches}

The same scope rule applies to the other downstream topics. Baryogenesis estimates built from powers of a phenomenological suppression factor, black-hole spectroscopy templates built from an additional discrete-horizon branch, and large-$N_{\mathrm{edge}}$ string/worldsheet reorganizations all require premises that are not established by the end of Theorem~\ref{thm:master}. In particular:
\begin{enumerate}[leftmargin=2em]
\item baryogenesis needs a concrete out-of-equilibrium mechanism and a justified defect-counting dynamics;
\item black-hole spectroscopy needs a discrete-horizon branch with its own spacing law and selection rules;
\item string/worldsheet claims need a controlled large-$N_{\mathrm{edge}}$ regime distinct from the physical color rank $N_c=3$.
\end{enumerate}
Until those extra premises are proved, these subjects remain program-level continuations only.

\subsection{Reading rule for the compact note}

The compact note should therefore be read with the following boundary in mind: relativity and Standard Model structure are the recovered core; the D6 capacity relation and the supplement-backed D10--D11 quantitative implementation are secondary; dark-sector, baryogenesis, spectroscopy, and string/worldsheet discussions are deferred continuations whose success or failure does not by itself confirm or refute the recovered core.

\section{Relation to Holography and Existing UV Frameworks}

OPH belongs to the holographic family of ideas, but its primitive data are different from those used in asymptotic-boundary constructions. The distinction matters because several of the central claims of the present paper, including the treatment of $\Lambda$, the handling of factorization, and the emergence of gauge structure, depend on it.

\subsection{Static-patch screen versus asymptotic boundary}

In asymptotic-boundary holography one starts from a dual theory at infinity and reconstructs the bulk inward. In OPH one starts from a finite-capacity screen equipped with a net of local patch algebras and reconstructs the effective bulk from overlap consistency. The primitive object is therefore a family of observer patches with nontrivial overlaps, not a single global boundary theory.

This shift has two technical consequences. First, subsystem factorization is treated from the beginning as a gluing problem with centers and sector labels. Second, the physical role of the horizon is local and operational: every observer has direct access only to a patch algebra, and global law is whatever survives reconciliation on overlaps.

\subsection{Positive $\Lambda$ and finite capacity}

The null-modular reconstruction of Section~\ref{sec:gravity} determines the stress tensor only up to a metric term. This means that local null data do not fix $\Lambda$. The cosmological constant enters only when one supplies the global screen capacity
\[
N_{\mathrm{scr}}=\log\dim\mathcal H_{\mathrm{tot}}.
\]
The resulting relation
\[
\Lambda=\frac{3\pi}{G N_{\mathrm{scr}}}
\]
fits naturally with a de Sitter static patch of finite entropy via $N_{\mathrm{scr}}=A_{\mathrm{dS}}/(4G)$. OPH is therefore formulated in a de Sitter-first language. The positive cosmological constant is tied to finite capacity rather than to a deformation of a negative-$\Lambda$ starting point.

\subsection{Factorization, gauge structure, and the string sector}

The factorization problem of gauge theory and gravity is often treated as a technical nuisance. In OPH it is part of the architecture. Edge-center completion turns the collar center into a first-class object, and the entropy split
\[
S_{\mathrm{gen}}=\langle L_C\rangle + S_{\mathrm{bulk}}
\]
is then a structural consequence rather than an added prescription.

In the bosonic EFT branch, the same collar data produce a symmetric rigid transportable sector category. Compact gauge reconstruction then follows conditionally from that category, and the Standard Model quotient is selected by minimal admissibility. This differs sharply from approaches in which the gauge group is specified in advance.

The string-theory relation is similarly reversed. The edge partition function
\[
Z_{\mathrm{edge}}(t)=\sum_R d_R^2 e^{-t C_2(R)}
\]
is the closed partition function obtained after gluing the open-edge weights $p_R(t)\propto d_R e^{-tC_2(R)}$. It has the heat-kernel form familiar from two-dimensional Yang-Mills theory. If a large-$N_{\mathrm{edge}}$ regime exists, with $N_{\mathrm{edge}}$ distinct from the physical color rank $N_c=3$, its free energy admits a genus expansion. Within OPH this is read as a conditional effective reorganization of edge degrees of freedom. String theory therefore enters as a sector description, while the microscopic starting data are the screen net, overlap consistency, and sector reconstruction.

\begingroup
\footnotesize
\begin{tabular}{@{}>{\raggedright\arraybackslash}p{0.24\textwidth}>{\raggedright\arraybackslash}p{0.32\textwidth}>{\raggedright\arraybackslash}p{0.34\textwidth}@{}}
\toprule
Feature & Asymptotic-boundary holography & OPH \\
\midrule
Primitive data & Boundary theory at infinity & Finite-capacity screen with observer-patch net \\
Factorization across cuts & Subtle and often indirect & Explicit edge-center completion with central labels \\
Status of $\Lambda$ & External background datum & Fixed globally by $N_{\mathrm{scr}}$ after local null reconstruction \\
Gauge group & Usually specified as part of the model & Reconstructed from edge sectors, then filtered by admissibility \\
String sector & Fundamental dual description & Effective large-$N_{\mathrm{edge}}$ reorganization of edge dynamics \\
\bottomrule
\end{tabular}
\endgroup

\section{Predictive Status and Falsifiability}

A scope-cleaned falsifiability statement has to track Theorem~\ref{thm:master} and the dependency ledger rather than the full visible list D1--D12. In this compact note, the recovered core is
\[
(D1\text{--}D5)\cup(D7\text{--}D9),
\]
namely the relativity chain together with the realized Standard Model structural chain. D6 is a separate input-dependent capacity corollary, D10--D11 are implementation-level numerical branches of the current supplement-backed realization, and D12 collects deferred continuations.

\subsection{Prediction audit by claim tier}

\begingroup
\footnotesize
\begin{longtable}{@{}>{\raggedright\arraybackslash}p{0.20\textwidth}>{\raggedright\arraybackslash}p{0.31\textwidth}>{\raggedright\arraybackslash}p{0.39\textwidth}@{}}
\toprule
Tier & Representative outputs & What would actually falsify the tier \\
\midrule
\endhead
\bottomrule
\endfoot
Phase I recovered core (Theorem~\ref{thm:master}; D1--D5, D7--D9) & Confluence of overlap repair, conditional Lorentz kinematics, the conditional Jacobson-type Einstein branch in the stated scaling regime, compact gauge reconstruction in the bosonic branch, the realized Standard Model quotient chain, the exact hypercharge lattice on the realized matter package, \(N_c=3\), \(N_g=3\), and the product-group consequence of no gauge-mediated proton decay & A mathematical failure in the derivation chain, or data that require a different realized gauge quotient, different realized hypercharges, a different color or generation count on the admissible branch, or gauge-mediated proton decay \\
Phase II input-dependent corollary (D6) & \(\Lambda=\frac{3\pi}{G N_{\mathrm{scr}}}\) once the screen-capacity identification is made, together with the separate order-of-magnitude neutrino estimate tied to the same capacity input & Failure of the screen-capacity identification or of the specific global capacity relation. Such a failure would not by itself erase the recovered core \\
Phase II current implementation: calibration sector (D10) & Pixel-closure gauge/electroweak consistency, \(\alpha_i(m_Z)\), \(\alpha_{\mathrm{em}}^{-1}(m_Z)\), \(\sin^2\theta_W(m_Z)\), \(v\), and tree-level electroweak masses & Failure of the supplement-backed closure of the present implementation once \(P\) and the printed running/matching conventions are imposed \\
Phase II current implementation: secondary quantitative branch (D11) & Higgs/top critical-surface branch & Failure of the extra critical-surface condition or of the supplement-backed RG transport used in that branch; such a failure does not by itself erase the recovered core \\
Phase III deferred continuations (D12 and beyond) & Flavor ans\"atze, charged-lepton/Koide fits, texture branches, neutrino mass/mixing constructions beyond the D6 estimate, dark-sector response laws, baryogenesis estimates, black-hole spectroscopy templates, and string/worldsheet reorganizations & Retraction of the corresponding continuation only. These branches are not part of the compact paper's recovered core and their failure is not, by itself, a falsification of relativity-plus-Standard-Model recovery \\
\end{longtable}
\endgroup

This tier split is the sync rule for mirrored status surfaces: Phase I is the recovered-core falsifiability standard, Phase II is visible but non-core, and Phase III is explicitly continuation-level.

\subsection{Recovered-core discriminators}

Once the claim boundary is made explicit, the sharpest discriminators are the ones that do not depend on auxiliary flavor, dark-sector, baryogenesis, or string ans\"atze.

\subsubsection{Realized Standard Model gauge structure}

On the realized MAR-admissible branch, the compact note fixes
\[
G_{\mathrm{phys}}=
\frac{\SU(3)\times\SU(2)\times\U(1)}{\mathbb Z_6},
\qquad
N_c=3,
\qquad
N_g=3,
\]
with the exact one-generation hypercharge lattice
\[
Y_Q=\frac16,\qquad
Y_L=-\frac12,\qquad
Y_u=-\frac23,\qquad
Y_d=\frac13,\qquad
Y_e=1,\qquad
Y_H=\frac12.
\]
Evidence that the realized low-energy gauge structure or realized hypercharges differ from these values would directly contradict the recovered core.

\subsubsection{Product-group corollary}

Because the realized gauge structure is a product group up to the finite central quotient already fixed above, there are no mixed \((3,2,\pm 5/6)\) gauge bosons. Hence
\[
\tau_p^{(\mathrm{gauge})}=\infty
\]
for gauge-mediated proton decay. Observation of proton decay specifically attributable to gauge-boson exchange would therefore falsify the realized product-group branch.

\subsubsection{Relativity branch}

The relativity claim is conditional but still sharp about its scope: if the stated geometric-modular, null-bridge, and fixed-cap-stationarity premises are realized in the intended scaling regime, then the compact note predicts Lorentz kinematics on the screen and a Jacobson-type Einstein branch. What is being tested here is the existence of that scaling regime and its premises, not a fixed-cutoff matrix-algebra identity.

\subsection{What does \emph{not} count as a contradiction of the recovered core}

The following do \emph{not} falsify Theorem~\ref{thm:master} if they fail in their present form:
\begin{enumerate}[leftmargin=2em]
\item the uniform $\mathbb Z_6$ center-label ensemble and the associated $\varepsilon=1/6$ flavor ansatz;
\item Koide/charged-lepton constructions, texture exponents, or capacity-level neutrino estimates;
\item modular-anomaly dark-sector response laws or MOND/RAR-style scaling claims;
\item baryogenesis estimates based on extra suppression-counting assumptions;
\item discrete-horizon spectroscopy templates;
\item large-$N_{\mathrm{edge}}$ string/worldsheet reorganizations.
\end{enumerate}
These are all post-Phase-I, non-core branches. Some are Phase-II input-dependent estimates, while the rest are Phase-III continuations. None is part of the compact note's recovered relativity-plus-Standard-Model core.

\subsection{What the compact paper claims after scope-cleaning}

After the present scope clean, the compact paper claims exactly this: observer-overlap consistency, together with the stated scaling-limit and categorical premises, recovers a conditional relativity branch and the realized Standard Model structural branch. The D6 capacity relation and the D10--D11 numerical implementation remain visible but secondary. Everything in Phase III is explicitly non-core.

\section{Common Objections and Clarifications}

\subsection{\texorpdfstring{Is the use of $P$ circular?}{Is the use of P circular?}}

The answer depends on which sector is under discussion. In the current implementation the pixel area $P$ calibrates the overall particle-physics scale through the gauge sector. Quantities that are algebraically entangled with that calibration, such as $v$, $m_W$, $m_Z$, and the gauge couplings at $m_Z$, are consistency checks. The paper treats them as such.

The stronger outputs enter only after the scale has been set. The critical-surface condition
\[
\lambda(M_U)=0,\qquad \beta_\lambda(M_U)=0
\]
fixes the Higgs and top branch without introducing a new continuous fit beyond the calibrated gauge trajectory and scale-setting branch. The Koide continuation is controlled by
\[
Q=\frac23,\qquad \delta=\frac29,
\]
which are continuation inputs tied to the previously derived discrete data inside that branch. The quark hierarchy depends on the $\mathbb Z_6$ texture ansatz, and the neutrino branch depends on the second input $N_{\mathrm{scr}}$, not on $P$.

Operationally, changing $P$ moves the calibration family $(\alpha_U,\alpha_i(m_Z),v,m_W,m_Z)$ as well as the overall GeV scale. It does not alter the parameter-free structural outputs such as the gauge quotient or hypercharge lattice. The useful distinction is therefore between scale setting, structural determination, and later phenomenological continuations.

\subsection{Does a fixed UV cell size break Lorentz invariance?}

The UV regulator is placed on the screen algebra, not on an emergent bulk spacetime lattice. Lorentz kinematics in the present framework arise from the continuum modular action on caps, summarized in Corollary~\ref{cor:lorentz}, rather than from microscopic invariance of a bulk discretization.

This is conceptually similar to other continuum limits in statistical mechanics and lattice field theory. A regulator may break a symmetry at finite cutoff while the infrared fixed point restores it exactly. In OPH the restoration is stronger than a generic RG expectation because the modular-flow theorem identifies the symmetry group itself:
\[
\Conf^+(S^2)\cong \mathrm{SO}^+(3,1).
\]
Residual Lorentz violation would therefore have to arise from a failure of the regulator premises or from a failure of the modular-flow theorem, not merely from the existence of a UV cell size.

\subsection{Why use type-I algebras if continuum QFT uses type III?}

The type-I assumption is a regulator statement. At finite cutoff the local patch algebras are finite-dimensional matrix algebras, which makes entropy and recovery theory well defined. The continuum limit is expected to recover the familiar type-III behavior of local QFT algebras. There is no contradiction here. The two statements refer to different scales.

This is no more problematic than using a lattice regulator for Yang-Mills theory and then taking a continuum limit in which the operator algebra changes type. The mathematical burden is the existence of a controlled scaling regime. In the present manuscript that burden is isolated in Assumption~\ref{ass:technical}. It is not hidden inside the main theorems.

\subsection{Why is charge quantization possible without a simple GUT?}

The framework uses the global quotient and anomaly structure rather than a simple-group embedding. Once the realized matter content is fixed, the subgroup acting trivially on all states is exactly $\mathbb Z_6$, and the anomaly equations force the sixth-integer hypercharge lattice. Integer electric charge for color singlets is then a property of the realized quotient structure. A simple unified gauge group is not required.

\section{Discussion}

The framework is strongest where the outputs are discrete or structurally rigid:
\begin{enumerate}[leftmargin=2em]
\item the Lorentz branch inside the assumed geometric modular phase;
\item the conditional scaling-limit Einstein branch under fixed-cap stationarity and the null-stress premise;
\item the Standard Model gauge quotient chain;
\item exact hypercharges on the realized one-generation matter package with one Higgs doublet;
\item three colors and three generations on the realized MAR-admissible branch;
\item the compact two-input economy of the present quantitative implementation, kept explicit as a secondary layer rather than as part of Phase I.
\end{enumerate}

The next strongest sector is the secondary quantitative D11 branch consisting of the Higgs and top masses. The separate D6 cosmological-capacity relation remains visible as an input-dependent corollary rather than as part of the recovered core. The charged-lepton and Koide sectors are promising but should presently be read as phenomenological continuations rather than core theorem outputs.

The weakest part is the light-quark and hadron branch. The hierarchy pattern is robust, but precision control of low-energy QCD effects, scheme dependence, and order-one coefficients is not yet complete.

The framework should therefore be evaluated in stages. It already contains several results that are normally postulated independently in other approaches: kinematics, dynamics, gauge structure, charge lattice, matter multiplicities, and part of the mass spectrum in supplement-backed quantitative form. Whether one regards this as evidence for fundamentality depends on whether these simultaneous derivations survive further scrutiny.

Topics intentionally excluded from the present formulation include simulation narratives, observer-continuation mechanisms, and claims of final completeness for all low-energy observables.\footnote{Questions outside this theorem package are, for present purposes, mostly harmless.}

\section{Conclusion}

Observer-Patch Holography starts from overlap consistency of observer patches and yields, within one framework, several effective structures usually treated separately in modern fundamental physics. The compact note supports a schedule-independent normal form for objective law, a conditional Lorentz branch in the assumed geometric modular phase, a conditional Jacobson-type Einstein branch under fixed-cap stationarity and the null-stress premise, the Standard Model gauge quotient with exact hypercharges on the realized matter package under the extended MAR admissibility package in the bosonic internal-gauge branch, three colors and three generations on that realized branch, a current two-input quantitative implementation whose supplement-backed Higgs/top branch is secondary to the recovered core, and a separate input-dependent cosmological-capacity corollary. It also records several phenomenological continuations whose present status is weaker and is labeled as such in the text.

Accordingly, the framework functions as a candidate underlying theory from which general relativity and the Standard Model may arise as effective descriptions once the stated scaling-limit and categorical premises are realized. The next technical tasks are sharper quantitative control, a fuller fermionic gauge branch, and explicit UV completions that land in the required modular/geometric phase.

\appendix

\section{One-loop RG Equations Used in the Quantitative Branch}\label{app:rges}

For convenience we record the one-loop system used in the Higgs/top transport branch:
Throughout this appendix, $g_Y$ denotes the SM hypercharge coupling. The GUT-normalized coupling is $g_1$ with $g_1^2=\frac53 g_Y^2$, equivalently $\alpha_1=\frac53\alpha_Y$. Therefore the one-loop coefficient is $41/6$ in the $g_Y$ equation and $41/10$ in the $g_1$ equation.
\begin{align}
16\pi^2\frac{dg_Y}{d\ln\mu} &= \frac{41}{6}g_Y^3,\\
16\pi^2\frac{dg_2}{d\ln\mu} &= -\frac{19}{6}g_2^3,\\
16\pi^2\frac{dg_3}{d\ln\mu} &= -7g_3^3,\\
16\pi^2\frac{dy_t}{d\ln\mu} &=
y_t\left(\frac92 y_t^2-\frac{17}{12}g_Y^2-\frac94 g_2^2-8g_3^2\right),\\
16\pi^2\frac{d\lambda}{d\ln\mu} &=
24\lambda^2+\lambda\left(12y_t^2-9g_2^2-3g_Y^2\right)-6y_t^4
+\frac38\left(2g_2^4+(g_2^2+g_Y^2)^2\right).
\end{align}
In the present supplement-backed numerical implementation this system is the transport layer between the UV boundary data and the low-energy Higgs/top outputs. Numerically, the compact note uses one-loop $\overline{\mathrm{MS}}$ transport, the matching scale $\mu=m_Z=91.1876\,\mathrm{GeV}$, the unified scale $M_U\approx 2.474\times 10^{16}\,\mathrm{GeV}$ from the main text, and the standard three-loop QCD pole-mass conversion for the displayed top-branch number. Higher-loop and threshold refinements are expected to shift the final numbers at the GeV level rather than alter the logic of the branch.

\section{Quantitative Robustness Audit}\label{app:robustness}

\begingroup
\footnotesize
\begin{tabular}{@{}>{\raggedright\arraybackslash}p{0.34\textwidth}>{\raggedright\arraybackslash}p{0.22\textwidth}>{\raggedright\arraybackslash}p{0.34\textwidth}@{}}
\toprule
Source of variation & Baseline used here & Expected effect on outputs \\
\midrule
Critical-surface RG transport & one-loop SM running & moving to two loops is expected to shift $m_H$ by about $1$ GeV and $m_t$ by about $0.5$ GeV \\
Pole-mass and threshold matching & standard three-loop top conversion plus leading electroweak corrections & order-one GeV effects in the top branch; sub-GeV to GeV effects in the Higgs branch \\
Heat-kernel truncation in $\bar\ell_G$ & 200 representations & negligible at displayed precision, with changes below $10^{-12}$ in the entropy sums \\
Bisection tolerance for the pixel closure & $10^{-10}$ & affects only the last shown digits of $\alpha_U$ and derived calibration outputs \\
Light-quark scheme matching & not yet implemented at precision level & dominant source of the present 5\% to 20\% uncertainty in the light-quark branch \\
\bottomrule
\end{tabular}
\endgroup

The paper therefore separates the quantitative sector by status. The Higgs and top branches are useful because their errors are controlled at the percent level. The charged-lepton continuation is numerically sharp but rests on a weaker discrete-selection step, while the light-quark and hadronic branch still carries the largest scheme and threshold uncertainty.

\begin{thebibliography}{99}

\bibitem{BW75}
J.~J.~Bisognano and E.~H.~Wichmann,
``On the duality condition for a Hermitian scalar field,''
\emph{J. Math. Phys.} \textbf{16} (1975) 985--1007.

\bibitem{BW76}
J.~J.~Bisognano and E.~H.~Wichmann,
``On the duality condition for quantum fields,''
\emph{J. Math. Phys.} \textbf{17} (1976) 303--321.

\bibitem{LiebRuskai}
E.~H.~Lieb and M.~B.~Ruskai,
``Proof of the strong subadditivity of quantum-mechanical entropy,''
\emph{J. Math. Phys.} \textbf{14} (1973) 1938--1941.

\bibitem{Petz86}
D.~Petz,
``Sufficient subalgebras and the relative entropy of states of a von Neumann algebra,''
\emph{Commun. Math. Phys.} \textbf{105} (1986) 123--131.

\bibitem{Petz88}
D.~Petz,
``Sufficiency of channels over von Neumann algebras,''
\emph{Quart. J. Math.} \textbf{39} (1988) 97--108.

\bibitem{FawziRenner}
O.~Fawzi and R.~Renner,
``Quantum conditional mutual information and approximate Markov chains,''
\emph{Commun. Math. Phys.} \textbf{340} (2015) 575--611.

\bibitem{HJPW}
P.~Hayden, R.~Jozsa, D.~Petz and A.~Winter,
``Structure of states which satisfy strong subadditivity of quantum entropy with equality,''
\emph{Commun. Math. Phys.} \textbf{246} (2004) 359--374.

\bibitem{Jacobson95}
T.~Jacobson,
``Thermodynamics of spacetime: the Einstein equation of state,''
\emph{Phys. Rev. Lett.} \textbf{75} (1995) 1260--1263.

\bibitem{Jacobson16}
T.~Jacobson,
``Entanglement equilibrium and the Einstein equation,''
\emph{Phys. Rev. Lett.} \textbf{116} (2016) 201101.

\bibitem{DR89}
S.~Doplicher and J.~E.~Roberts,
``A new duality theory for compact groups,''
\emph{Invent. Math.} \textbf{98} (1989) 157--218.

\bibitem{DR90}
S.~Doplicher and J.~E.~Roberts,
``Why there is a field algebra with a compact gauge group describing the superselection structure in particle physics,''
\emph{Commun. Math. Phys.} \textbf{131} (1990) 51--107.

\bibitem{Witten82}
E.~Witten,
``An SU(2) anomaly,''
\emph{Phys. Lett. B} \textbf{117} (1982) 324--328.

\end{thebibliography}

\end{document}
